\documentclass[11pt]{article}

\usepackage[a4paper,margin=1in]{geometry}
\usepackage{amsmath,amssymb,amsthm,mathtools}
\usepackage[T1]{fontenc}
\usepackage{lmodern}
\usepackage{microtype}
\usepackage{hyperref}

\newtheorem{theorem}{Theorem}
\newtheorem{lemma}[theorem]{Lemma}
\newtheorem{proposition}[theorem]{Proposition}
\newtheorem{corollary}[theorem]{Corollary}
\newtheorem{remark}[theorem]{Remark}
\newtheorem{definition}[theorem]{Definition}
\newtheorem{hypothesis}[theorem]{Hypothesis}
\newtheorem{conjecture}[theorem]{Conjecture}

\title{Topological Problems for the Clean Hatano--Nelson Family with Open Boundary Conditions}
\author{}
\date{}

\begin{document}
\maketitle

For \(n>1\), \(0<a<1\), consider the real nonnormal tridiagonal Toeplitz matrix
\[
\mathbf{A}_n(a)\coloneqq \operatorname{tridiag}(a,0,1)
=
\begin{bmatrix}
0 & 1 \\
a & 0 & \ddots \\
& \ddots & \ddots & \ddots \\
&& \ddots & 0 & 1 \\
&&& a & 0
\end{bmatrix}\in \mathbb{R}^{n\times n}.
\]
For \(\varepsilon>0\), its \(\varepsilon\)-pseudospectrum is
\[
\sigma_\varepsilon\bigl(\mathbf{A}_n(a)\bigr)
\coloneqq
\left\{
z\in\mathbb C:\;
\bigl\|(z\mathbf I_n-\mathbf A_n(a))^{-1}\bigr\|_2>\varepsilon^{-1}
\right\},
\]
with the usual convention that the resolvent norm is \(+\infty\) on the spectrum. Equivalently,
\[
\sigma_\varepsilon\bigl(\mathbf A_n(a)\bigr)
=
\left\{
z\in\mathbb C:\;
s_{\min}\!\bigl(z\mathbf I_n-\mathbf A_n(a)\bigr)<\varepsilon
\right\}.
\]

Throughout, for \(z=x+iy\) we write
\[
\Sigma_\varepsilon\coloneqq \sigma_\varepsilon\bigl(\mathbf A_n(a)\bigr),
\qquad
s(x,y)\coloneqq s_{\min}\!\bigl((x+iy)\mathbf I_n-\mathbf A_n(a)\bigr).
\]
Then
\[
\Sigma_\varepsilon=\{x+iy\in\mathbb C:\; s(x,y)<\varepsilon\}.
\]

The original simply-connectedness statement is reduced below to a vertical monotonicity problem. The topological implication is rigorous; the remaining monotonicity statement is formulated as a conjecture and then rewritten in several equivalent ways.

\section{Topological reduction to vertical monotonicity}

\begin{conjecture}[Vertical monotonicity]\label{conj:vertical-monotonicity}
For every fixed \(x\in\mathbb R\), the function
\[
[0,\infty)\ni y\longmapsto s(x,y)
\]
is nondecreasing.
\end{conjecture}

\begin{theorem}\label{thm:components-simply-connected}
Assume Conjecture~\ref{conj:vertical-monotonicity}. Then for every \(\varepsilon>0\), every connected component of \(\Sigma_\varepsilon\) is simply connected.
\end{theorem}

We first record several rigorous lemmas.

\begin{lemma}\label{lem:real-spectrum}
With
\[
\mathbf D=\operatorname{diag}\bigl(1,\sqrt a,a,\dots,a^{(n-1)/2}\bigr),
\]
one has
\[
\mathbf D^{-1}\mathbf A_n(a)\mathbf D
=
\sqrt a\,\operatorname{tridiag}(1,0,1).
\]
Hence the eigenvalues of \(\mathbf A_n(a)\) are real and simple, namely
\[
\lambda_k
=
2\sqrt a\cos\!\left(\frac{k\pi}{n+1}\right),
\qquad
k=1,\dots,n.
\]
\end{lemma}

\begin{proof}
A direct computation gives
\[
(\mathbf D^{-1}\mathbf A_n(a)\mathbf D)_{j,j+1}
=
a^{-(j-1)/2}a^{j/2}
=
\sqrt a,
\]
and
\[
(\mathbf D^{-1}\mathbf A_n(a)\mathbf D)_{j+1,j}
=
a^{-j/2}a\,a^{(j-1)/2}
=
\sqrt a.
\]
Thus \(\mathbf D^{-1}\mathbf A_n(a)\mathbf D\) is the symmetric tridiagonal Toeplitz matrix \(\sqrt a\,\operatorname{tridiag}(1,0,1)\). Its eigenvalues are the classical Dirichlet eigenvalues \(2\sqrt a\cos(k\pi/(n+1))\), and they are simple.
\end{proof}

\begin{lemma}\label{lem:real-axis}
Every connected component of \(\Sigma_\varepsilon\) meets the real axis.
\end{lemma}

\begin{proof}
Let \(\Omega\) be a connected component of \(\Sigma_\varepsilon\). Since
\[
s_{\min}\!\bigl(z\mathbf I_n-\mathbf A_n(a)\bigr)\ge |z|-\|\mathbf A_n(a)\|_2,
\]
the set \(\Sigma_\varepsilon\) is bounded; hence \(\Omega\) is bounded.

Assume that \(\Omega\cap\sigma(\mathbf A_n(a))=\varnothing\). Define
\[
\varphi(z)\coloneqq \bigl\|(z\mathbf I_n-\mathbf A_n(a))^{-1}\bigr\|_2.
\]
Then \(\varphi\) is subharmonic on \(\Omega\) and continuous on \(\overline\Omega\). Moreover,
\[
\varphi(z)>\varepsilon^{-1}
\qquad (z\in\Omega).
\]

We claim that \(\varphi(z)=\varepsilon^{-1}\) on \(\partial\Omega\). Indeed, \(\Omega\) is a connected component of the open set \(\Sigma_\varepsilon\), hence \(\partial\Omega\subset \partial\Sigma_\varepsilon\). Since \(s\) is continuous and
\[
\Sigma_\varepsilon=\{s<\varepsilon\},
\]
one has \(s=\varepsilon\) on \(\partial\Sigma_\varepsilon\), therefore on \(\partial\Omega\). Since \(\Omega\cap\sigma(\mathbf A_n(a))=\varnothing\), no point of \(\partial\Omega\) can belong to the spectrum: an eigenvalue would satisfy \(s=0<\varepsilon\) and hence would lie in \(\Sigma_\varepsilon\), whereas \(\partial\Omega\subset \mathbb C\setminus \Omega\). Thus
\[
\varphi(z)=\varepsilon^{-1}
\qquad (z\in\partial\Omega).
\]

This contradicts the maximum principle for subharmonic functions on bounded domains. Therefore \(\Omega\) must contain an eigenvalue of \(\mathbf A_n(a)\). By Lemma~\ref{lem:real-spectrum}, every eigenvalue is real, so
\[
\Omega\cap\mathbb R\neq\varnothing.
\]
\end{proof}

\begin{lemma}\label{lem:conjugation-symmetry}
For every \(x,y\in\mathbb R\),
\[
s(x,-y)=s(x,y).
\]
Consequently \(\Sigma_\varepsilon\) is invariant under complex conjugation.
\end{lemma}

\begin{proof}
Because \(\mathbf A_n(a)\) has real entries,
\[
\overline{(x+iy)\mathbf I_n-\mathbf A_n(a)}
=
(x-iy)\mathbf I_n-\mathbf A_n(a).
\]
Entrywise complex conjugation preserves singular values, hence
\[
s(x,-y)
=
s_{\min}\!\bigl((x-iy)\mathbf I_n-\mathbf A_n(a)\bigr)
=
s_{\min}\!\bigl((x+iy)\mathbf I_n-\mathbf A_n(a)\bigr)
=
s(x,y).
\]
\end{proof}

\begin{lemma}\label{lem:vertical-sections}
Assume Conjecture~\ref{conj:vertical-monotonicity}. For every \(x\in\mathbb R\), the vertical section
\[
V_x
\coloneqq
\{y\in\mathbb R:\; x+iy\in\Sigma_\varepsilon\}
\]
is either empty or an open interval symmetric about \(0\).
\end{lemma}

\begin{proof}
By Lemma~\ref{lem:conjugation-symmetry}, \(V_x\) is symmetric under \(y\mapsto -y\). By Conjecture~\ref{conj:vertical-monotonicity}, the set
\[
\{y\ge 0:\; s(x,y)<\varepsilon\}
\]
is either empty or an interval of the form \([0,h_x)\) with \(h_x\in(0,\infty]\). Since \(\Sigma_\varepsilon\) is bounded, in fact \(h_x<\infty\). Therefore
\[
V_x=\varnothing
\quad\text{or}\quad
V_x=(-h_x,h_x).
\]
\end{proof}

\begin{lemma}\label{lem:component-fibers}
Assume Conjecture~\ref{conj:vertical-monotonicity}. Let \(\Omega\) be a connected component of \(\Sigma_\varepsilon\). Then there is an open interval \(I_\Omega\subset\mathbb R\) and a function
\[
h_\Omega:I_\Omega\to(0,\infty)
\]
such that
\[
\Omega
=
\{x+iy:\; x\in I_\Omega,\ |y|<h_\Omega(x)\}.
\]
\end{lemma}

\begin{proof}
Define
\[
I_\Omega
\coloneqq
\{x\in\mathbb R:\; (\{x\}+i\mathbb R)\cap\Omega\neq\varnothing\}.
\]
Since the projection map \(x+iy\mapsto x\) is continuous and \(\Omega\) is connected, \(I_\Omega\) is connected. Since \(\Omega\) is open, \(I_\Omega\) is open: if \(x_0\in I_\Omega\), then \(x_0+iy_0\in\Omega\) for some \(y_0\), and an open disk about \(x_0+iy_0\) contained in \(\Omega\) projects onto an open interval about \(x_0\). Hence \(I_\Omega\) is an open interval.

By Lemma~\ref{lem:real-axis}, there exists \(x_0\in\Omega\cap\mathbb R\). Complex conjugation is a homeomorphism of \(\Sigma_\varepsilon\), so it sends \(\Omega\) onto another connected component of \(\Sigma_\varepsilon\). Because \(x_0\) is fixed by conjugation, the image component contains \(x_0\), hence must coincide with \(\Omega\). Therefore \(\Omega\) is invariant under complex conjugation.

For each \(x\in I_\Omega\), let
\[
\Omega_x
\coloneqq
\{y\in\mathbb R:\; x+iy\in\Omega\}.
\]
Then \(\Omega_x\) is a nonempty open subset of \(V_x\), and it is symmetric under \(y\mapsto -y\).

By Lemma~\ref{lem:vertical-sections}, \(V_x\) is either empty or a connected open interval. Moreover, \(V_x\) is partitioned by the disjoint open sets
\[
\Omega'_x\coloneqq \{y\in\mathbb R:\; x+iy\in\Omega'\},
\]
where \(\Omega'\) ranges over the connected components of \(\Sigma_\varepsilon\). Since \(V_x\) is connected, at most one component \(\Omega'\) can meet the vertical line \(\{x\}+i\mathbb R\). Because \(\Omega\) does meet that line, it follows that
\[
\Omega_x=V_x.
\]

Therefore, for each \(x\in I_\Omega\), the fiber \(\Omega_x\) is an interval symmetric about \(0\). Write
\[
\Omega_x=(-h_\Omega(x),h_\Omega(x)),
\qquad h_\Omega(x)>0.
\]
Then
\[
\Omega
=
\{x+iy:\; x\in I_\Omega,\ |y|<h_\Omega(x)\}.
\]
\end{proof}

\begin{proof}[Proof of Theorem~\ref{thm:components-simply-connected}]
Let \(\Omega\) be a connected component of \(\Sigma_\varepsilon\). By Lemma~\ref{lem:component-fibers},
\[
\Omega
=
\{x+iy:\; x\in I_\Omega,\ |y|<h_\Omega(x)\},
\]
where \(I_\Omega\subset\mathbb R\) is an open interval.

Define
\[
H:[0,1]\times\Omega\longrightarrow\Omega,
\qquad
H(t,x+iy)=x+i(1-t)y.
\]
For fixed \(x\in I_\Omega\), the vertical fiber of \(\Omega\) above \(x\) is \((-h_\Omega(x),h_\Omega(x))\), so \(H(t,x+iy)\) stays in that same fiber for every \(t\in[0,1]\). Therefore \(H\) is a deformation retraction of \(\Omega\) onto
\[
\Omega\cap\mathbb R=I_\Omega.
\]
Hence \(\Omega\) is contractible. In particular, \(\Omega\) is simply connected.
\end{proof}

\section{Singular-vector identities on simple branches}

Let
\[
\mathbf R\in\mathbb R^{n\times n},
\qquad
\mathbf R_{j,n+1-j}=1,
\qquad
\mathbf R_{jk}=0\ \text{otherwise},
\]
be the reversal matrix.

\begin{lemma}\label{lem:persymmetry}
One has
\[
\mathbf R\,\mathbf A_n(a)\,\mathbf R=\mathbf A_n(a)^{\mathsf T}.
\]
Suppose \(x\in\mathbb R\) is fixed and \(I\subset(0,\infty)\) is an interval on which
\[
s(x,y)=s_{\min}\!\bigl((x+iy)\mathbf I_n-\mathbf A_n(a)\bigr)
\]
is simple. Then the associated unit left and right singular vectors \(\mathbf u(y),\mathbf v(y)\) may be chosen differentiably on \(I\) so that
\[
\bigl((x+iy)\mathbf I_n-\mathbf A_n(a)\bigr)\mathbf v(y)=s(x,y)\,\mathbf u(y),
\]
\[
\bigl((x+iy)\mathbf I_n-\mathbf A_n(a)\bigr)^*\mathbf u(y)=s(x,y)\,\mathbf v(y),
\]
and
\[
\mathbf u(y)=\mathbf R\,\overline{\mathbf v(y)}.
\]
\end{lemma}

\begin{proof}
The identity \(\mathbf R\,\mathbf A_n(a)\,\mathbf R=\mathbf A_n(a)^{\mathsf T}\) is immediate from the Toeplitz form of \(\mathbf A_n(a)\).

On an interval of simplicity, one may choose differentiable unit eigenvectors of the Hermitian matrix
\[
\mathbf B_y^*\mathbf B_y,
\qquad
\mathbf B_y\coloneqq (x+iy)\mathbf I_n-\mathbf A_n(a),
\]
corresponding to the simple eigenvalue \(s(x,y)^2\); denote such a vector by \(\mathbf v(y)\). Since \(y>0\), Lemma~\ref{lem:real-spectrum} implies \(s(x,y)>0\), and we may define
\[
\mathbf u(y)\coloneqq s(x,y)^{-1}\mathbf B_y\mathbf v(y).
\]
Then \(\mathbf u(y)\) is differentiable, has unit norm, and satisfies
\[
\mathbf B_y\mathbf v=s\,\mathbf u,
\qquad
\mathbf B_y^*\mathbf u=s\,\mathbf v,
\]
with \(s=s(x,y)\).

Since \(\mathbf R\mathbf B_y=\mathbf B_y^{\mathsf T}\mathbf R\), we get
\[
\mathbf B_y(\mathbf R\overline{\mathbf u})
=
\mathbf R\,\mathbf B_y^{\mathsf T}\overline{\mathbf u}
=
\mathbf R\,\overline{\mathbf B_y^*\mathbf u}
=
s\,\mathbf R\overline{\mathbf v},
\]
and similarly
\[
\mathbf B_y^*(\mathbf R\overline{\mathbf v})=s\,\mathbf R\overline{\mathbf u}.
\]
Thus \((\mathbf R\overline{\mathbf u},\mathbf R\overline{\mathbf v})\) is another singular pair for the same simple singular value \(s\). By simplicity, there is \(\alpha(y)\in\mathbb C\), \(|\alpha(y)|=1\), such that
\[
\mathbf R\overline{\mathbf u}=\alpha\,\mathbf v,
\qquad
\mathbf R\overline{\mathbf v}=\alpha\,\mathbf u.
\]
Choose a differentiable branch \(\beta(y)\) of a square root of \(\alpha(y)\), and replace \((\mathbf u,\mathbf v)\) by \((\beta\mathbf u,\beta\mathbf v)\). Then
\[
\mathbf R\overline{\beta\mathbf v}
=
\overline{\beta}\,\mathbf R\overline{\mathbf v}
=
\overline{\beta}\,\alpha\,\mathbf u
=
\beta\,\mathbf u.
\]
Hence, after this common phase normalization,
\[
\mathbf u=\mathbf R\,\overline{\mathbf v}.
\]
\end{proof}

\begin{proposition}[Derivative formula and current identity]\label{prop:current-identity}
Fix \(x\in\mathbb R\), let \(I\subset(0,\infty)\) be an interval on which \(s(x,y)\) is simple, and choose \(\mathbf u(y),\mathbf v(y)\) as in Lemma~\ref{lem:persymmetry}. Write
\[
\mathbf v(y)=
\begin{bmatrix}
v_1(y)\\
\vdots\\
v_n(y)
\end{bmatrix},
\]
and define
\[
J_j(y)\coloneqq \Im\!\bigl(\overline{v_j(y)}\,v_{j+1}(y)\bigr),
\qquad
j=1,\dots,n-1,
\]
with \(J_0(y)=J_n(y)=0\). Then:
\begin{enumerate}
\item
\[
\frac{d}{dy}s(x,y)
=
-\Im\!\bigl(\mathbf u(y)^*\mathbf v(y)\bigr)
=
-\Im\!\bigl(\mathbf v(y)^{\mathsf T}\mathbf R\,\mathbf v(y)\bigr).
\]
\item For each \(j=1,\dots,n\),
\[
J_j-aJ_{j-1}
=
y|v_j|^2+s(x,y)\,\Im\!\bigl(v_jv_{n+1-j}\bigr).
\]
\item Summing over \(j\) gives
\[
(1-a)\sum_{j=1}^{n-1}J_j
=
y+s(x,y)\,\Im\!\bigl(\mathbf u(y)^*\mathbf v(y)\bigr).
\]
Consequently,
\[
\frac{d}{dy}\bigl(s(x,y)^2\bigr)
=
2y-2(1-a)\sum_{j=1}^{n-1}J_j.
\]
\end{enumerate}
\end{proposition}

\begin{proof}
Let
\[
\mathbf B_y\coloneqq (x+iy)\mathbf I_n-\mathbf A_n(a).
\]
Since \(\mathbf B_y' = i\mathbf I_n\), the standard singular-value perturbation formula gives
\[
\frac{d}{dy}s(x,y)
=
\Re\!\bigl(\mathbf u^* \mathbf B_y' \mathbf v\bigr)
=
\Re\!\bigl(i\,\mathbf u^*\mathbf v\bigr)
=
-\Im\!\bigl(\mathbf u^*\mathbf v\bigr).
\]
By Lemma~\ref{lem:persymmetry}, \(\mathbf u=\mathbf R\overline{\mathbf v}\), so
\[
\mathbf u^*\mathbf v
=
\mathbf v^{\mathsf T}\mathbf R\,\mathbf v.
\]
This proves (1).

For (2), the singular-vector equation \(\mathbf B_y\mathbf v=s\,\mathbf u\) reads componentwise
\[
(x+iy)v_j-a v_{j-1}-v_{j+1}
=
s\,\overline{v_{n+1-j}},
\qquad
j=1,\dots,n,
\]
with Dirichlet endpoints \(v_0=v_{n+1}=0\). Multiplying the \(j\)-th equation by \(\overline{v_j}\) and taking imaginary parts yields
\[
y|v_j|^2-a\,\Im(\overline{v_j}v_{j-1})-\Im(\overline{v_j}v_{j+1})
=
-s\,\Im(v_jv_{n+1-j}).
\]
Since
\[
\Im(\overline{v_j}v_{j-1})=-\Im(\overline{v_{j-1}}v_j)=-J_{j-1},
\qquad
\Im(\overline{v_j}v_{j+1})=J_j,
\]
we obtain
\[
J_j-aJ_{j-1}
=
y|v_j|^2+s\,\Im(v_jv_{n+1-j}).
\]
This proves (2).

Summing (2) over \(j=1,\dots,n\) and using \(J_0=J_n=0\) and \(\|\mathbf v\|_2=1\), we get
\[
(1-a)\sum_{j=1}^{n-1}J_j
=
y+s\,\sum_{j=1}^n \Im(v_jv_{n+1-j}).
\]
But
\[
\sum_{j=1}^n v_jv_{n+1-j}
=
\mathbf v^{\mathsf T}\mathbf R\,\mathbf v
=
\mathbf u^*\mathbf v,
\]
hence
\[
(1-a)\sum_{j=1}^{n-1}J_j
=
y+s\,\Im(\mathbf u^*\mathbf v).
\]
Combining this with (1) gives
\[
\frac{d}{dy}s(x,y)
=
\frac{y-(1-a)\sum_{j=1}^{n-1}J_j}{s(x,y)},
\]
and therefore
\[
\frac{d}{dy}\bigl(s(x,y)^2\bigr)
=
2s(x,y)\frac{d}{dy}s(x,y)
=
2y-2(1-a)\sum_{j=1}^{n-1}J_j.
\]
\end{proof}

\section{Exact reformulations of vertical monotonicity}

Let \(\mathbf S\in\mathbb C^{n\times n}\) be the upper shift,
\[
\mathbf S_{j,j+1}=1,
\qquad
\mathbf S_{jk}=0\ \text{otherwise}.
\]
Then
\[
\mathbf A_n(a)=a\mathbf S^{\mathsf T}+\mathbf S.
\]
We also define
\[
\mathbf K\coloneqq \frac{1-a}{2i}\,(\mathbf S-\mathbf S^{\mathsf T}),
\qquad
\mathbf M_x\coloneqq (x\mathbf I_n-\mathbf A_n(a))^*(x\mathbf I_n-\mathbf A_n(a)).
\]

\begin{proposition}\label{prop:H-affine}
For every \(x,y\in\mathbb R\),
\[
\mathbf H_x(y)\coloneqq \bigl((x+iy)\mathbf I_n-\mathbf A_n(a)\bigr)^*\bigl((x+iy)\mathbf I_n-\mathbf A_n(a)\bigr)
=
\mathbf M_x+y^2\mathbf I_n-2y\,\mathbf K.
\]
In particular,
\[
s(x,y)^2=\lambda_{\min}\!\bigl(\mathbf H_x(y)\bigr).
\]
\end{proposition}

\begin{proof}
Since
\[
\bigl((x+iy)\mathbf I_n-\mathbf A_n(a)\bigr)^*
=
(x-iy)\mathbf I_n-\mathbf A_n(a)^{\mathsf T},
\]
we have
\[
\mathbf H_x(y)
=
\mathbf M_x
+y^2\mathbf I_n
+iy\bigl(\mathbf A_n(a)-\mathbf A_n(a)^{\mathsf T}\bigr).
\]
Now
\[
\mathbf A_n(a)-\mathbf A_n(a)^{\mathsf T}
=
(1-a)(\mathbf S-\mathbf S^{\mathsf T}),
\]
hence
\[
iy\bigl(\mathbf A_n(a)-\mathbf A_n(a)^{\mathsf T}\bigr)
=
-2y\,\frac{1-a}{2i}(\mathbf S-\mathbf S^{\mathsf T})
=
-2y\,\mathbf K.
\]
Therefore
\[
\mathbf H_x(y)=\mathbf M_x+y^2\mathbf I_n-2y\,\mathbf K.
\]
Since \(\mathbf H_x(y)=\mathbf B_y^*\mathbf B_y\), its eigenvalues are the squares of the singular values of \(\mathbf B_y=(x+iy)\mathbf I_n-\mathbf A_n(a)\), and the smallest one is \(s(x,y)^2\).
\end{proof}

\begin{corollary}[Scalar form of vertical monotonicity]\label{cor:scalar-form}
Let \(x\in\mathbb R\), and suppose \(s(x,y)^2\) is simple on an interval \(I\subset(0,\infty)\). Let \(\boldsymbol{\psi}(y)\) be a normalized eigenvector of \(\mathbf H_x(y)\) for
\[
\lambda(y)=s(x,y)^2.
\]
Then
\[
\lambda'(y)=2y-2\langle \boldsymbol{\psi}(y),\mathbf K\boldsymbol{\psi}(y)\rangle.
\]
Consequently, on such an interval the vertical monotonicity statement
\[
\frac{d}{dy}s(x,y)\ge 0
\]
is equivalent to
\[
\langle \boldsymbol{\psi}(y),\mathbf K\boldsymbol{\psi}(y)\rangle\le y.
\]
If \(\boldsymbol{\psi}(y)=\mathbf v(y)\) is the right singular vector from Proposition~\ref{prop:current-identity}, then
\[
\langle \boldsymbol{\psi}(y),\mathbf K\boldsymbol{\psi}(y)\rangle
=
(1-a)\sum_{j=1}^{n-1}\Im\!\bigl(\overline{v_j(y)}\,v_{j+1}(y)\bigr).
\]
\end{corollary}

\begin{proof}
By Proposition~\ref{prop:H-affine},
\[
\mathbf H_x'(y)=2y\,\mathbf I_n-2\mathbf K.
\]
The Hellmann--Feynman formula therefore gives
\[
\lambda'(y)=\langle \boldsymbol{\psi},\mathbf H_x'(y)\boldsymbol{\psi}\rangle
=
2y-2\langle \boldsymbol{\psi},\mathbf K\boldsymbol{\psi}\rangle.
\]
Since \(\lambda=s^2\) and \(s>0\) on \(y>0\), one has
\[
\frac{d}{dy}s(x,y)\ge 0
\quad\Longleftrightarrow\quad
\frac{d}{dy}\bigl(s(x,y)^2\bigr)\ge 0
\quad\Longleftrightarrow\quad
\langle \boldsymbol{\psi},\mathbf K\boldsymbol{\psi}\rangle\le y.
\]

Finally, if \(\boldsymbol{\psi}=\mathbf v\), then
\[
\langle \mathbf v,\mathbf K\mathbf v\rangle
=
\frac{1-a}{2i}\,\mathbf v^*(\mathbf S-\mathbf S^{\mathsf T})\mathbf v
=
\frac{1-a}{2i}\sum_{j=1}^{n-1}\bigl(\overline{v_j}v_{j+1}-\overline{v_{j+1}}v_j\bigr)
=
(1-a)\sum_{j=1}^{n-1}\Im(\overline{v_j}v_{j+1}).
\]
\end{proof}

\begin{remark}[Concavity after subtracting \(y^2\)]\label{rem:concavity}
By Proposition~\ref{prop:H-affine},
\[
s(x,y)^2-y^2
=
\lambda_{\min}(\mathbf M_x-2y\,\mathbf K)
=
\min_{\|\boldsymbol{\xi}\|_2=1}\bigl(\langle \boldsymbol{\xi},\mathbf M_x\boldsymbol{\xi}\rangle-2y\langle \boldsymbol{\xi},\mathbf K\boldsymbol{\xi}\rangle\bigr).
\]
Thus \(y\mapsto s(x,y)^2-y^2\) is the pointwise infimum of affine functions of \(y\), hence is concave on \(\mathbb R\).
\end{remark}

\subsection*{A matrix-valued compression}

\begin{proposition}[Hidden matrix positivity on the minimal singular subspace]\label{prop:hidden-matrix-positivity}
Assume \(s=s(x,y)>0\), and let \(m\) be the multiplicity of the singular value \(s\) of
\[
\mathbf B_y\coloneqq (x+iy)\mathbf I_n-\mathbf A_n(a).
\]
Let
\[
\mathbf V\in\mathbb C^{n\times m}
\]
have orthonormal columns spanning the right singular subspace
\[
\ker(\mathbf B_y^*\mathbf B_y-s^2\mathbf I_n),
\]
and define
\[
\mathbf U\coloneqq s^{-1}\mathbf B_y\mathbf V.
\]
Then \(\mathbf U^*\mathbf U=\mathbf V^*\mathbf V=\mathbf I_m\), and
\[
\mathbf B_y\mathbf V=s\,\mathbf U,
\qquad
\mathbf B_y^*\mathbf U=s\,\mathbf V.
\]
Define the Hermitian linearization
\[
\mathcal B_x(y)\coloneqq
\begin{bmatrix}
0 & \mathbf B_y\\
\mathbf B_y^* & 0
\end{bmatrix},
\qquad
\mathcal J\coloneqq
\begin{bmatrix}
0 & i\mathbf I_n\\
-i\mathbf I_n & 0
\end{bmatrix},
\]
and
\[
\mathbf W\coloneqq \frac{1}{\sqrt2}\begin{bmatrix}\mathbf U\\[1mm]\mathbf V\end{bmatrix}\in\mathbb C^{2n\times m}.
\]
Then
\[
\mathcal B_x(y)\mathbf W=s\,\mathbf W,
\qquad
\mathbf W^*\mathbf W=\mathbf I_m,
\]
and
\[
\boxed{
\mathbf W^*\mathcal J\mathbf W
=
\frac1s
\left(
y\,\mathbf I_m
-
\frac{1-a}{2i}\,\mathbf V^*(\mathbf S-\mathbf S^{\mathsf T})\mathbf V
\right).
}
\]
In particular, when \(m=1\), this reduces exactly to the scalar inequality from Corollary~\ref{cor:scalar-form}.
\end{proposition}

\begin{proof}
The relations \(\mathbf U^*\mathbf U=\mathbf I_m\), \(\mathbf B_y\mathbf V=s\,\mathbf U\), and \(\mathbf B_y^*\mathbf U=s\,\mathbf V\) are immediate. Therefore
\[
\mathcal B_x(y)\mathbf W
=
\frac{1}{\sqrt2}
\begin{bmatrix}
\mathbf B_y\mathbf V\\
\mathbf B_y^*\mathbf U
\end{bmatrix}
=
\frac{1}{\sqrt2}
\begin{bmatrix}
s\,\mathbf U\\
s\,\mathbf V
\end{bmatrix}
=
s\,\mathbf W,
\]
and
\[
\mathbf W^*\mathbf W
=
\frac12(\mathbf U^*\mathbf U+\mathbf V^*\mathbf V)
=
\mathbf I_m.
\]

Next,
\[
\mathbf W^*\mathcal J\mathbf W
=
\frac12
\begin{bmatrix}\mathbf U^* & \mathbf V^*\end{bmatrix}
\begin{bmatrix}
0 & i\mathbf I_n\\
-i\mathbf I_n & 0
\end{bmatrix}
\begin{bmatrix}\mathbf U\\ \mathbf V\end{bmatrix}
=
\frac{i}{2}\,(\mathbf U^*\mathbf V-\mathbf V^*\mathbf U).
\]
Since \(\mathbf U=s^{-1}\mathbf B_y\mathbf V\),
\[
\mathbf U^*\mathbf V-\mathbf V^*\mathbf U
=
s^{-1}\mathbf V^*(\mathbf B_y^*-\mathbf B_y)\mathbf V.
\]
Now
\[
\mathbf B_y^*-\mathbf B_y
=
-2iy\,\mathbf I_n+\mathbf A_n(a)-\mathbf A_n(a)^{\mathsf T}
=
-2iy\,\mathbf I_n+(1-a)(\mathbf S-\mathbf S^{\mathsf T}).
\]
Hence
\[
\mathbf W^*\mathcal J\mathbf W
=
\frac{i}{2s}\,
\mathbf V^*\bigl(-2iy\,\mathbf I_n+(1-a)(\mathbf S-\mathbf S^{\mathsf T})\bigr)\mathbf V
=
\frac1s
\left(
y\,\mathbf I_m
-
\frac{1-a}{2i}\,\mathbf V^*(\mathbf S-\mathbf S^{\mathsf T})\mathbf V
\right).
\]
\end{proof}

\begin{remark}[Quadratic pencil positivity]\label{rem:quadratic-pencil}
For fixed \(x\in\mathbb R\) and \(\varepsilon>0\), define the Hermitian quadratic pencil
\[
\mathbf Q_{x,\varepsilon}(y)
\coloneqq
\mathbf H_x(y)-\varepsilon^2\mathbf I_n
=
\mathbf M_x+y^2\mathbf I_n-2y\,\mathbf K-\varepsilon^2\mathbf I_n.
\]
Then
\[
x+iy\in\partial\Sigma_\varepsilon
\quad\Longleftrightarrow\quad
\mathbf Q_{x,\varepsilon}(y)\succeq 0
\ \text{ and }\
\ker \mathbf Q_{x,\varepsilon}(y)\neq\{0\}.
\]
Indeed, \(\partial\Sigma_\varepsilon\) is exactly the level set \(s(x,y)=\varepsilon\), and
\[
s(x,y)=\varepsilon
\quad\Longleftrightarrow\quad
\lambda_{\min}(\mathbf H_x(y))=\varepsilon^2
\quad\Longleftrightarrow\quad
\mathbf Q_{x,\varepsilon}(y)\succeq 0,\ \ker\mathbf Q_{x,\varepsilon}(y)\neq\{0\}.
\]

If \(0\) is a simple eigenvalue of \(\mathbf Q_{x,\varepsilon}(y)\) with normalized eigenvector \(\boldsymbol{\psi}\), then
\[
\langle \boldsymbol{\psi},\mathbf Q_{x,\varepsilon}'(y)\boldsymbol{\psi}\rangle
=
2\bigl(y-\langle \boldsymbol{\psi},\mathbf K\boldsymbol{\psi}\rangle\bigr).
\]
Thus the desired vertical monotonicity at a boundary point is equivalent to the nonnegativity of the compression of \(\mathbf Q_{x,\varepsilon}'(y)\) to \(\ker \mathbf Q_{x,\varepsilon}(y)\). This is the quadratic-pencil form of the hidden positivity.
\end{remark}

\section{Rigorous partial results}

\begin{proposition}[A large-\(y\) monotonicity region]\label{prop:large-y}
Let \(\lambda(y)=s(x,y)^2\), and let \(\boldsymbol{\psi}(y)\) be a normalized eigenvector of \(\mathbf H_x(y)\) for \(\lambda(y)\) on an interval of simplicity. Then
\[
\lambda'(y)\ge 0
\qquad\text{whenever}\qquad
y\ge (1-a)\cos\!\left(\frac{\pi}{n+1}\right).
\]
\end{proposition}

\begin{proof}
By Corollary~\ref{cor:scalar-form},
\[
\lambda'(y)=2y-2\langle \boldsymbol{\psi},\mathbf K\boldsymbol{\psi}\rangle.
\]
Since \(\mathbf K\) is Hermitian,
\[
\langle \boldsymbol{\psi},\mathbf K\boldsymbol{\psi}\rangle\le \|\mathbf K\|.
\]
Now
\[
\mathbf K=\frac{1-a}{2}\,\mathbf L,
\qquad
\mathbf L\coloneqq i(\mathbf S^{\mathsf T}-\mathbf S).
\]
Let
\[
\mathbf D=\operatorname{diag}\bigl(1,i,i^2,\dots,i^{n-1}\bigr).
\]
A direct computation shows
\[
\mathbf D^*\mathbf L\mathbf D=\mathbf S+\mathbf S^{\mathsf T}.
\]
Thus
\[
\|\mathbf K\|
=
\frac{1-a}{2}\,\|\mathbf S+\mathbf S^{\mathsf T}\|.
\]
The symmetric tridiagonal Toeplitz matrix \(\mathbf S+\mathbf S^{\mathsf T}\) has eigenvalues
\[
2\cos\!\left(\frac{k\pi}{n+1}\right),\qquad k=1,\dots,n,
\]
hence
\[
\|\mathbf S+\mathbf S^{\mathsf T}\|=2\cos\!\left(\frac{\pi}{n+1}\right).
\]
Therefore
\[
\|\mathbf K\|
=
(1-a)\cos\!\left(\frac{\pi}{n+1}\right).
\]
If \(y\ge \|\mathbf K\|\), then
\[
\lambda'(y)=2y-2\langle \boldsymbol{\psi},\mathbf K\boldsymbol{\psi}\rangle\ge 2y-2\|\mathbf K\|\ge 0.
\]
\end{proof}

\begin{proposition}[The case \(n=2\)]\label{prop:n=2}
For \(n=2\), one has
\[
\frac{d}{dy}\bigl(s(x,y)^2\bigr)\ge 0
\qquad (y\ge 0,\ x\in\mathbb R).
\]
In particular, Conjecture~\ref{conj:vertical-monotonicity} holds for \(n=2\).
\end{proposition}

\begin{proof}
For \(n=2\),
\[
\mathbf B_y=
\begin{bmatrix}
x+iy & -1\\
-a & x+iy
\end{bmatrix},
\]
so
\[
\mathbf H_x(y)=\mathbf B_y^*\mathbf B_y=
\begin{bmatrix}
a^2+x^2+y^2 & -(a+1)x+i(1-a)y\\
-(a+1)x-i(1-a)y & 1+x^2+y^2
\end{bmatrix}.
\]
The smaller eigenvalue is
\[
s(x,y)^2
=
x^2+y^2+\frac{1+a^2}{2}
-\frac12\sqrt{(1-a^2)^2+4(a+1)^2x^2+4(1-a)^2y^2}.
\]
Differentiating gives
\[
\frac{d}{dy}\bigl(s(x,y)^2\bigr)
=
2y\left(
1-\frac{(1-a)^2}{
\sqrt{(1-a^2)^2+4(a+1)^2x^2+4(1-a)^2y^2}}
\right).
\]
Since
\[
\sqrt{(1-a^2)^2+4(a+1)^2x^2+4(1-a)^2y^2}\ge 1-a^2\ge (1-a)^2,
\]
the bracket is nonnegative, hence
\[
\frac{d}{dy}\bigl(s(x,y)^2\bigr)\ge 0
\qquad (y\ge 0).
\]
Because \(s(x,y)\ge 0\), this implies \(d s(x,y)/dy\ge 0\) whenever \(s(x,y)>0\), and therefore Conjecture~\ref{conj:vertical-monotonicity} follows for \(n=2\).
\end{proof}

\section{Convexity and the reduced-resolvent form}

Fix \(x\in\mathbb R\). On an interval of simplicity, let
\[
\lambda(y)\coloneqq s(x,y)^2=\lambda_{\min}(\mathbf H_x(y)),
\]
and let \(\boldsymbol{\psi}(y)\) be a normalized eigenvector of \(\mathbf H_x(y)\) for \(\lambda(y)\). Write the remaining eigenpairs as
\[
\mathbf H_x(y)\boldsymbol{\psi}_k(y)=\lambda_k(y)\boldsymbol{\psi}_k(y),
\qquad
k=2,\dots,n,
\]
with \(\{\boldsymbol{\psi},\boldsymbol{\psi}_2,\dots,\boldsymbol{\psi}_n\}\) orthonormal, and define the positive reduced resolvent
\[
\mathcal G(y)\coloneqq \sum_{k=2}^n \frac{1}{\lambda_k(y)-\lambda(y)}\,\boldsymbol{\psi}_k(y)\boldsymbol{\psi}_k(y)^*.
\]
Equivalently,
\[
(\mathbf H_x(y)-\lambda(y)\mathbf I_n)\mathcal G(y)=\mathcal G(y)(\mathbf H_x(y)-\lambda(y)\mathbf I_n)=\mathbf I_n-\boldsymbol{\psi}(y)\boldsymbol{\psi}(y)^*.
\]

\begin{proposition}[Reduced-resolvent form of convexity]\label{prop:reduced-resolvent}
On an interval of simplicity,
\[
\lambda''(y)
=
2-8\,\bigl\langle \mathbf K\boldsymbol{\psi}(y),\,\mathcal G(y)\,\mathbf K\boldsymbol{\psi}(y)\bigr\rangle.
\]
Hence the convexity statement
\[
\lambda''(y)\ge 0
\]
is equivalent to
\[
\boxed{
\bigl\langle \mathbf K\boldsymbol{\psi}(y),\,\mathcal G(y)\,\mathbf K\boldsymbol{\psi}(y)\bigr\rangle\le \frac14.
}
\]
\end{proposition}

\begin{proof}
We have
\[
\mathbf H_x'(y)=2y\,\mathbf I_n-2\mathbf K,
\qquad
\mathbf H_x''(y)=2\mathbf I_n.
\]
For a simple Hermitian eigenvalue branch, the second-derivative formula is
\[
\lambda''(y)
=
\langle \boldsymbol{\psi},\mathbf H_x''(y)\boldsymbol{\psi}\rangle
-
2\sum_{k=2}^n \frac{\bigl|\langle \boldsymbol{\psi}_k,\mathbf H_x'(y)\boldsymbol{\psi}\rangle\bigr|^2}{\lambda_k-\lambda}.
\]
Since \(\boldsymbol{\psi}_k^*\boldsymbol{\psi}=0\), the scalar part \(2y\mathbf I_n\) drops out:
\[
\langle \boldsymbol{\psi}_k,\mathbf H_x'(y)\boldsymbol{\psi}\rangle
=
-2\langle \boldsymbol{\psi}_k,\mathbf K\boldsymbol{\psi}\rangle.
\]
Therefore
\[
\lambda''(y)
=
2-8\sum_{k=2}^n \frac{|\langle \boldsymbol{\psi}_k,\mathbf K\boldsymbol{\psi}\rangle|^2}{\lambda_k-\lambda}
=
2-8\,\langle \mathbf K\boldsymbol{\psi},\mathcal G\,\mathbf K\boldsymbol{\psi}\rangle.
\]
This is exactly the stated bound.
\end{proof}

\begin{remark}[A numerical-range reformulation]\label{rem:numerical-range}
Define
\[
\alpha_x(t)\coloneqq
\min\bigl\{
\langle \boldsymbol{\xi},\mathbf M_x\boldsymbol{\xi}\rangle:\;
\|\boldsymbol{\xi}\|_2=1,\ \langle \boldsymbol{\xi},\mathbf K\boldsymbol{\xi}\rangle=t
\bigr\},
\]
for \(t\) in the interval
\[
I_x\coloneqq \{\langle \boldsymbol{\xi},\mathbf K\boldsymbol{\xi}\rangle:\ \|\boldsymbol{\xi}\|_2=1\}.
\]
Then
\[
\lambda(y)-y^2
=
\min_{t\in I_x}\bigl(\alpha_x(t)-2yt\bigr).
\]
If \(y\mapsto \lambda(y)\) is simple and the lower boundary function \(\alpha_x\) is \(C^2\) at the minimizing point \(t(y)=\langle \boldsymbol{\psi}(y),\mathbf K\boldsymbol{\psi}(y)\rangle\), then
\[
\alpha_x'(t(y))=2y,
\qquad
\lambda''(y)-2=-\frac{4}{\alpha_x''(t(y))}.
\]
Comparing with Proposition~\ref{prop:reduced-resolvent}, one obtains
\[
\bigl\langle \mathbf K\boldsymbol{\psi}(y),\,\mathcal G(y)\,\mathbf K\boldsymbol{\psi}(y)\bigr\rangle
=
\frac{1}{2\,\alpha_x''(t(y))}.
\]
Thus convexity of \(\lambda(y)=s(x,y)^2\) is equivalent to the lower boundary of this numerical-range profile being \(2\)-strongly convex:
\[
\alpha_x''(t)\ge 2.
\]
\end{remark}

\section{Boundary-defect and boundary-resolvent reductions}

Define
\[
\mathbf T\coloneqq \mathbf S+\mathbf S^{\mathsf T},
\qquad
\mathbf L\coloneqq i(\mathbf S^{\mathsf T}-\mathbf S),
\qquad
\mathbf P_1\coloneqq \mathbf e_1\mathbf e_1^{\mathsf T},
\qquad
\mathbf P_n\coloneqq \mathbf e_n\mathbf e_n^{\mathsf T}.
\]
Then
\[
\mathbf K=\frac{1-a}{2}\,\mathbf L,
\qquad
\mathbf A_n(a)=\frac{1+a}{2}\,\mathbf T+i\mathbf K.
\]

\begin{proposition}[Boundary-defect identity]\label{prop:boundary-defect}
For every \(x,y\in\mathbb R\),
\[
\mathbf H_x(y)
=
\Bigl(x\mathbf I_n-\frac{1+a}{2}\mathbf T\Bigr)^2
+
\bigl(y\mathbf I_n-\mathbf K\bigr)^2
+
\frac{1-a^2}{2}\,(\mathbf P_n-\mathbf P_1).
\]
\end{proposition}

\begin{proof}
Write
\[
\mathbf X\coloneqq x\mathbf I_n-\frac{1+a}{2}\mathbf T,
\qquad
\mathbf Y\coloneqq y\mathbf I_n-\mathbf K.
\]
Then
\[
(x+iy)\mathbf I_n-\mathbf A_n(a)=\mathbf X+i\mathbf Y.
\]
Hence
\[
\mathbf H_x(y)
=
(\mathbf X-i\mathbf Y)(\mathbf X+i\mathbf Y)
=
\mathbf X^2+\mathbf Y^2+i[\mathbf X,\mathbf Y].
\]
Since scalar matrices commute,
\[
[\mathbf X,\mathbf Y]
=
\left[-\frac{1+a}{2}\mathbf T,\,-\mathbf K\right]
=
\frac{1+a}{2}\,[\mathbf T,\mathbf K].
\]
Now
\[
\mathbf K=\frac{1-a}{2}\mathbf L,
\]
so
\[
[\mathbf T,\mathbf K]=\frac{1-a}{2}[\mathbf T,\mathbf L].
\]
A direct computation using
\[
\mathbf S\mathbf S^{\mathsf T}=\mathbf I_n-\mathbf P_n,
\qquad
\mathbf S^{\mathsf T}\mathbf S=\mathbf I_n-\mathbf P_1,
\]
gives
\[
[\mathbf T,\mathbf L]
=
2i(\mathbf P_1-\mathbf P_n).
\]
Therefore
\[
i[\mathbf X,\mathbf Y]
=
i\cdot \frac{1+a}{2}\cdot \frac{1-a}{2}\cdot 2i(\mathbf P_1-\mathbf P_n)
=
\frac{1-a^2}{2}\,(\mathbf P_n-\mathbf P_1).
\]
Substituting this into the expression for \(\mathbf H_x(y)\) proves the claim.
\end{proof}

\begin{remark}
Proposition~\ref{prop:boundary-defect} shows that the noncommutative correction to the ``sum of squares''
\[
\Bigl(x\mathbf I_n-\frac{1+a}{2}\mathbf T\Bigr)^2
+
\bigl(y\mathbf I_n-\mathbf K\bigr)^2
\]
is a rank-two boundary defect supported at the left and right endpoints. This strongly suggests that the unresolved convexity issue is a boundary phenomenon rather than a bulk one.
\end{remark}

We next isolate an almost-commuting symmetry of
\[
\mathbf B_y=(x+iy)\mathbf I_n-\mathbf A_n(a).
\]

\begin{proposition}[Rank-two commutator for \(\mathbf B_y\)]\label{prop:B-commutator}
Let
\[
\mathbf M\coloneqq \mathbf P_n-\mathbf P_1.
\]
Then
\[
[\mathbf B_y,\mathbf L]=i(1+a)\mathbf M,
\qquad
[\mathbf B_y^*,\mathbf L]=i(1+a)\mathbf M.
\]
Consequently,
\[
[\mathbf H_x(y),\mathbf L]
=
i(1+a)\bigl(\mathbf B_y^*\mathbf M+\mathbf M\mathbf B_y\bigr).
\]
Moreover, the Hermitian matrix
\[
\mathbf C_x\coloneqq \mathbf B_y^*\mathbf M+\mathbf M\mathbf B_y
\]
is independent of \(y\) and is supported on \(\operatorname{span}\{\mathbf e_1,\mathbf e_2,\mathbf e_{n-1},\mathbf e_n\}\). Explicitly,
\[
\mathbf C_x
=
-2x\,\mathbf P_1
+2x\,\mathbf P_n
+\mathbf E_{12}+\mathbf E_{21}
-a\mathbf E_{n-1,n}
-a\mathbf E_{n,n-1},
\]
where \(\mathbf E_{jk}=\mathbf e_j\mathbf e_k^{\mathsf T}\).
\end{proposition}

\begin{proof}
Since \(\mathbf B_y=(x+iy)\mathbf I_n-a\mathbf S^{\mathsf T}-\mathbf S\), it is enough to compute
\[
[\mathbf S,\mathbf L]
=
\mathbf S\,i(\mathbf S^{\mathsf T}-\mathbf S)-i(\mathbf S^{\mathsf T}-\mathbf S)\mathbf S
=
i(\mathbf S\mathbf S^{\mathsf T}-\mathbf S^{\mathsf T}\mathbf S)
=
i(\mathbf P_1-\mathbf P_n),
\]
and similarly
\[
[\mathbf S^{\mathsf T},\mathbf L]
=
i(\mathbf P_1-\mathbf P_n).
\]
Hence
\[
[\mathbf B_y,\mathbf L]
=
-a[\mathbf S^{\mathsf T},\mathbf L]-[\mathbf S,\mathbf L]
=
-(1+a)i(\mathbf P_1-\mathbf P_n)
=
i(1+a)\mathbf M.
\]
The identity for \([\mathbf B_y^*,\mathbf L]\) is identical.

Now
\[
[\mathbf H_x(y),\mathbf L]
=
\mathbf B_y^*[\mathbf B_y,\mathbf L]+[\mathbf B_y^*,\mathbf L]\mathbf B_y
=
i(1+a)\bigl(\mathbf B_y^*\mathbf M+\mathbf M\mathbf B_y\bigr).
\]
This proves the commutator identity.

To compute \(\mathbf C_x\), note that
\[
\mathbf M\mathbf B_y
=
-x\,\mathbf P_1
+x\,\mathbf P_n
+\mathbf E_{12}
-a\mathbf E_{n,n-1},
\]
and
\[
\mathbf B_y^*\mathbf M
=
-x\,\mathbf P_1
+x\,\mathbf P_n
+\mathbf E_{21}
-a\mathbf E_{n-1,n}.
\]
The \(iy\)-terms cancel, so \(\mathbf C_x\) is independent of \(y\), and summing the two displays gives the stated formula.
\end{proof}

\begin{proposition}[Boundary Green-function reduction]\label{prop:boundary-green}
Let \(x\in\mathbb R\), and suppose \(\lambda(y)=s(x,y)^2\) is simple on an interval of \(y\)-values. Let \(\boldsymbol{\psi}(y)\) be a normalized eigenvector of \(\mathbf H_x(y)\) for \(\lambda(y)\), and let
\[
\mathcal G(y)
=
\sum_{k=2}^n \frac{1}{\lambda_k(y)-\lambda(y)}\,\boldsymbol{\psi}_k(y)\boldsymbol{\psi}_k(y)^*
\]
be the positive reduced resolvent. Then
\[
(\mathbf I_n-\boldsymbol{\psi}\boldsymbol{\psi}^*)\,\mathbf L\boldsymbol{\psi}
=
i(1+a)\,\mathcal G\,\mathbf C_x\boldsymbol{\psi},
\]
and
\[
\boxed{
\bigl\langle \mathbf K\boldsymbol{\psi},\mathcal G\,\mathbf K\boldsymbol{\psi}\bigr\rangle
=
\frac{(1-a^2)^2}{4}\,
\bigl\langle \mathbf C_x\boldsymbol{\psi},\mathcal G^{\,3}\mathbf C_x\boldsymbol{\psi}\bigr\rangle.
}
\]
In particular, by Proposition~\ref{prop:reduced-resolvent}, convexity of \(y\mapsto s(x,y)^2\) is equivalent to the boundary estimate
\[
(1-a^2)^2\,
\bigl\langle \mathbf C_x\boldsymbol{\psi},\mathcal G^{\,3}\mathbf C_x\boldsymbol{\psi}\bigr\rangle
\le 1.
\]
\end{proposition}

\begin{proof}
Since \(\mathbf H_x(y)\boldsymbol{\psi}=\lambda(y)\boldsymbol{\psi}\), Proposition~\ref{prop:B-commutator} gives
\[
(\mathbf H_x(y)-\lambda(y)\mathbf I_n)\,\mathbf L\boldsymbol{\psi}
=
[\mathbf H_x(y),\mathbf L]\boldsymbol{\psi}
=
i(1+a)\mathbf C_x\boldsymbol{\psi}.
\]
Now \(\mathbf C_x\) is Hermitian, and
\[
\langle \boldsymbol{\psi},[\mathbf H_x(y),\mathbf L]\boldsymbol{\psi}\rangle=0,
\]
so
\[
\langle \boldsymbol{\psi},\mathbf C_x\boldsymbol{\psi}\rangle=0.
\]
Hence \(\mathbf C_x\boldsymbol{\psi}\in \boldsymbol{\psi}^\perp\), and applying \(\mathcal G\) to the previous identity yields
\[
(\mathbf I_n-\boldsymbol{\psi}\boldsymbol{\psi}^*)\,\mathbf L\boldsymbol{\psi}
=
i(1+a)\,\mathcal G\,\mathbf C_x\boldsymbol{\psi}.
\]
Since
\[
\mathbf K=\frac{1-a}{2}\mathbf L,
\]
we obtain
\[
(\mathbf I_n-\boldsymbol{\psi}\boldsymbol{\psi}^*)\,\mathbf K\boldsymbol{\psi}
=
i\,\frac{1-a^2}{2}\,\mathcal G\,\mathbf C_x\boldsymbol{\psi}.
\]
Because \(\mathcal G\boldsymbol{\psi}=0\),
\[
\langle \mathbf K\boldsymbol{\psi},\mathcal G\,\mathbf K\boldsymbol{\psi}\rangle
=
\left\langle
(\mathbf I_n-\boldsymbol{\psi}\boldsymbol{\psi}^*)\mathbf K\boldsymbol{\psi},\,
\mathcal G\,(\mathbf I_n-\boldsymbol{\psi}\boldsymbol{\psi}^*)\mathbf K\boldsymbol{\psi}
\right\rangle.
\]
Substituting the previous display gives
\[
\langle \mathbf K\boldsymbol{\psi},\mathcal G\,\mathbf K\boldsymbol{\psi}\rangle
=
\frac{(1-a^2)^2}{4}\,
\langle \mathcal G\mathbf C_x\boldsymbol{\psi},\mathcal G^{\,2}\mathbf C_x\boldsymbol{\psi}\rangle
=
\frac{(1-a^2)^2}{4}\,
\langle \mathbf C_x\boldsymbol{\psi},\mathcal G^{\,3}\mathbf C_x\boldsymbol{\psi}\rangle.
\]
The last statement follows by combining this identity with Proposition~\ref{prop:reduced-resolvent}.
\end{proof}

\begin{remark}[Antiunitary symmetry of the ground state]\label{rem:antiunitary}
Since \(\mathbf H_x(y)\) has real coefficients except for purely imaginary skew-symmetric terms, one checks that
\[
\mathbf R\,\overline{\mathbf H_x(y)}\,\mathbf R=\mathbf H_x(y).
\]
Hence, whenever \(\lambda(y)\) is simple, the normalized eigenvector \(\boldsymbol{\psi}(y)\) may be chosen so that
\[
\mathbf R\,\overline{\boldsymbol{\psi}(y)}=\boldsymbol{\psi}(y).
\]
Equivalently,
\[
\psi_n=\overline{\psi_1},
\qquad
\psi_{n-1}=\overline{\psi_2}.
\]
With the explicit formula for \(\mathbf C_x\) from Proposition~\ref{prop:B-commutator},
\[
\mathbf C_x\boldsymbol{\psi}
=
(-2x\psi_1+\psi_2)\mathbf e_1
+\psi_1\mathbf e_2
-a\overline{\psi_1}\,\mathbf e_{n-1}
+\bigl(2x\overline{\psi_1}-a\overline{\psi_2}\bigr)\mathbf e_n.
\]
Thus the boundary source \(\mathbf C_x\boldsymbol{\psi}\) is determined by the first two components of \(\boldsymbol{\psi}\).
\end{remark}

\section{Constant-coefficient bulk recurrence}

\begin{proposition}[Interior fourth-order recurrence]\label{prop:bulk-recurrence}
Let \(\mu\in\mathbb R\), and let
\[
\mathbf u=
\begin{bmatrix}
u_1\\
\vdots\\
u_n
\end{bmatrix}
\]
satisfy
\[
(\mathbf H_x(y)-\mu\mathbf I_n)\mathbf u=0.
\]
Then for every interior index \(j\) for which all shifted terms are defined,
\[
a u_{j+2}
-c\,u_{j+1}
+
\bigl(1+a^2+x^2+y^2-\mu\bigr)u_j
-\overline c\,u_{j-1}
+a u_{j-2}
=0,
\]
where
\[
c\coloneqq (1+a)x-i(1-a)y.
\]
Equivalently, the characteristic polynomial of the bulk recurrence is
\[
p_\mu(r)
=
a r^4-c r^3+\bigl(1+a^2+x^2+y^2-\mu\bigr)r^2-\overline c\,r+a.
\]
It factorizes as
\[
\boxed{
p_\mu(r)
=
(ar^2-\overline z\,r+1)(r^2-zr+a)-\mu r^2,
\qquad
z=x+iy.
}
\]
\end{proposition}

\begin{proof}
Let
\[
\mathbf B_y=(x+iy)\mathbf I_n-\mathbf A_n(a).
\]
Then
\[
(\mathbf B_y\mathbf u)_j
=
zu_j-a u_{j-1}-u_{j+1},
\qquad z=x+iy.
\]
Since
\[
\mathbf H_x(y)=\mathbf B_y^*\mathbf B_y,
\]
one computes
\begin{align*}
(\mathbf H_x(y)\mathbf u)_j
&=
\overline z(zu_j-a u_{j-1}-u_{j+1})
-(zu_{j-1}-a u_{j-2}-u_j)
-a(zu_{j+1}-a u_j-u_{j+2}) \\
&=
a u_{j+2}
-(\overline z+a z)u_{j+1}
+\bigl(1+a^2+|z|^2\bigr)u_j
-(z+a\overline z)u_{j-1}
+a u_{j-2}.
\end{align*}
Now
\[
\overline z+a z=(1+a)x-i(1-a)y=c,
\qquad
z+a\overline z=\overline c.
\]
Thus
\[
(\mathbf H_x(y)-\mu\mathbf I_n)\mathbf u=0
\]
gives the stated recurrence. The polynomial \(p_\mu(r)\) is its characteristic polynomial. Expanding
\[
(ar^2-\overline z\,r+1)(r^2-zr+a)-\mu r^2
\]
shows that it is exactly equal to \(p_\mu(r)\).
\end{proof}

\begin{remark}
Proposition~\ref{prop:bulk-recurrence} shows that the bulk equation is a constant-coefficient fourth-order recurrence with a palindromic characteristic polynomial. The unresolved boundary-resolvent estimate from Proposition~\ref{prop:boundary-green} should therefore be expressible in terms of the two characteristic roots of \(p_\mu\) inside the unit disk and an explicit boundary Schur complement.
\end{remark}

\section{Empirical observations and conjectural next steps}

\begin{remark}[Empirical observations]\label{rem:numerics}
Informal numerical experiments suggest the following statements for all \(n>1\), \(0<a<1\), and fixed \(x\in\mathbb R\):
\begin{enumerate}
\item the map \(y\mapsto s(x,y)^2\) is convex on \([0,\infty)\);
\item the quantity
\[
\mathcal R_x(y)\coloneqq \bigl\langle \mathbf K\boldsymbol{\psi}(y),\,\mathcal G(y)\,\mathbf K\boldsymbol{\psi}(y)\bigr\rangle
\]
appears to be nonincreasing in \(y\);
\item on simple singular-value branches, the local currents
\[
J_j(y)=\Im\!\bigl(\overline{v_j(y)}\,v_{j+1}(y)\bigr)
\]
appear to satisfy \(J_j(y)\ge 0\), and the reflected phases
\[
\Im\!\bigl(v_j(y)v_{n+1-j}(y)\bigr)
\]
appear to satisfy
\[
\Im\!\bigl(v_j(y)v_{n+1-j}(y)\bigr)\le 0.
\]
\end{enumerate}
At present these statements remain conjectural. The rigorous reductions above show that the weakest statement needed for the topology is the scalar inequality
\[
\langle \boldsymbol{\psi}(y),\mathbf K\boldsymbol{\psi}(y)\rangle\le y,
\]
or equivalently the reduced-resolvent bound
\[
\bigl\langle \mathbf K\boldsymbol{\psi}(y),\,\mathcal G(y)\,\mathbf K\boldsymbol{\psi}(y)\bigr\rangle\le \frac14.
\]
\end{remark}

\begin{remark}[Summary of the remaining analytic problem]
The simply-connectedness theorem is reduced to Conjecture~\ref{conj:vertical-monotonicity}. On intervals of simplicity, this conjecture is equivalent to any of the following:
\begin{enumerate}
\item
\[
\langle \boldsymbol{\psi}(y),\mathbf K\boldsymbol{\psi}(y)\rangle\le y;
\]
\item
\[
\mathbf W^*\mathcal J\mathbf W\succeq 0
\quad\text{on the minimal positive eigenspace of the Hermitian linearization }\mathcal B_x(y);
\]
\item
\[
\mathbf Q'_{x,\varepsilon}(y)\succeq 0
\quad\text{on }\ker \mathbf Q_{x,\varepsilon}(y)
\quad\text{at boundary points }s(x,y)=\varepsilon;
\]
\item
\[
\lambda''(y)\ge 0
\quad\Longleftrightarrow\quad
\bigl\langle \mathbf K\boldsymbol{\psi}(y),\mathcal G(y)\mathbf K\boldsymbol{\psi}(y)\bigr\rangle\le \frac14;
\]
\item
\[
(1-a^2)^2\,
\bigl\langle \mathbf C_x\boldsymbol{\psi}(y),\mathcal G(y)^3\mathbf C_x\boldsymbol{\psi}(y)\bigr\rangle\le 1.
\]
\end{enumerate}
The last form is a \(4\times4\) boundary Green-function problem, because \(\mathbf C_x\) is supported only on \(\operatorname{span}\{\mathbf e_1,\mathbf e_2,\mathbf e_{n-1},\mathbf e_n\}\).
\end{remark}

\section{Exact boundary Weyl matrix from the root basis}

In this section we assume \(n\ge 5\).

Fix \(x\in\mathbb R\), \(y>0\), and write
\[
z=x+iy,
\qquad
\mathbf H_{x,y}
\coloneqq
\bigl(z\mathbf I_n-\mathbf A_n(a)\bigr)^*\bigl(z\mathbf I_n-\mathbf A_n(a)\bigr).
\]
Then
\[
\lambda_1(x,y)=s(x,y)^2=\lambda_{\min}(\mathbf H_{x,y}).
\]
Also set
\[
c\coloneqq \overline z + a z = (1+a)x-i(1-a)y,
\qquad
d_\mu\coloneqq 1+a^2+|z|^2-\mu.
\]
For \(\mu\in\mathbb R\), recall from Proposition~\ref{prop:bulk-recurrence} the quartic
\[
p_\mu(r)
=
a r^4-c r^3+d_\mu r^2-\overline c\,r+a.
\]

We introduce the boundary-coordinate map
\[
\mathbf P_\partial:\mathbb C^n\to\mathbb C^4,
\qquad
\mathbf P_\partial \mathbf u=
\begin{bmatrix}
u_1\\
u_2\\
u_{n-1}\\
u_n
\end{bmatrix}.
\]

\begin{proposition}[Where the two-decaying-root picture breaks]\label{prop:symbol-floor}
For \(\theta\in\mathbb R\), one has
\[
p_\mu(e^{i\theta})=0
\quad\Longleftrightarrow\quad
\mu=
\bigl|z-e^{i\theta}-a e^{-i\theta}\bigr|^2.
\]
Hence the quartic \(p_\mu\) has a unit-circle root if and only if
\[
\mu\in
\left\{
\bigl|z-e^{i\theta}-a e^{-i\theta}\bigr|^2:\ \theta\in[0,2\pi)
\right\}.
\]
In particular, if we define the symbol floor
\[
\mu_{\mathrm{sym}}(x,y)
\coloneqq
\min_{\theta\in[0,2\pi)}
\bigl|z-e^{i\theta}-a e^{-i\theta}\bigr|^2,
\]
then the regime in which \(p_\mu\) has exactly two roots in the open unit disk and two outside it is, in general, only
\[
\mu<\mu_{\mathrm{sym}}(x,y).
\]
\end{proposition}

\begin{proof}
For \(|r|=1\), one has \(r^{-1}=\overline r\), and therefore
\[
r^{-2}p_\mu(r)
=
a r^2-c r+d_\mu-\overline c\,r^{-1}+a r^{-2}.
\]
Using \(c=\overline z + a z\), a direct expansion gives
\[
a r^2-c r+(1+a^2+|z|^2)-\overline c\,r^{-1}+a r^{-2}
=
\bigl|z-r-a r^{-1}\bigr|^2.
\]
Hence, for \(|r|=1\),
\[
r^{-2}p_\mu(r)=\bigl|z-r-a r^{-1}\bigr|^2-\mu.
\]
Taking \(r=e^{i\theta}\) gives the claim.

Finally, \(p_\mu\) is self-inversive, with leading and constant coefficients both equal to \(a\), so its roots come in reciprocal-conjugate pairs and their product has modulus \(1\). Therefore, once there is no unit-circle root, there must be exactly two roots in \(\mathbb D\) and two outside \(\overline{\mathbb D}\).
\end{proof}

\begin{remark}
Proposition~\ref{prop:symbol-floor} shows that the correct exact boundary object is the \emph{four-root} Evans/Weyl matrix below, not the two-decaying-root formula by itself. The latter survives as an exact factorization only in the subthreshold regime \(\mu<\mu_{\mathrm{sym}}(x,y)\).
\end{remark}

\subsection*{The four-root boundary trace/residual matrices}

Assume for the moment that \(p_\mu\) has four simple roots
\[
\zeta_1(\mu),\dots,\zeta_4(\mu).
\]
(Repeated-root parameters are exceptional; the formulas below extend to them by continuity, or by replacing repeated root columns with the corresponding generalized-root solutions.)

For each root \(\zeta\) of \(p_\mu\), define the scalar sequence
\[
u_j^{(\zeta)}\coloneqq \zeta^{\,j-1},
\qquad
j=1,\dots,n.
\]
Because \(p_\mu(\zeta)=0\), the vector
\[
\mathbf u^{(\zeta)}=
\begin{bmatrix}
u_1^{(\zeta)}\\
\vdots\\
u_n^{(\zeta)}
\end{bmatrix}
\]
solves the interior homogeneous equation
\[
(\mathbf H_{x,y}-\mu\mathbf I_n)\mathbf u^{(\zeta)}=0
\qquad\text{for rows }j=3,\dots,n-2.
\]

Define the boundary trace column
\[
\mathbf t(\zeta)\coloneqq
\begin{bmatrix}
1\\
\zeta\\
\zeta^{\,n-2}\\
\zeta^{\,n-1}
\end{bmatrix}
=
\mathbf P_\partial \mathbf u^{(\zeta)}.
\]

Define also the boundary residual column
\[
\mathbf n(\zeta)\coloneqq
\mathbf P_\partial
\bigl((\mathbf H_{x,y}-\mu\mathbf I_n)\mathbf u^{(\zeta)}\bigr).
\]
A direct row-by-row computation gives
\[
\mathbf n(\zeta)=
\begin{bmatrix}
\overline c\,\zeta^{-1}-a\zeta^{-2}-1\\[1mm]
-a\zeta^{-1}\\[1mm]
-a\zeta^{\,n}\\[1mm]
\zeta^{\,n-1}\bigl(c\zeta-a\zeta^2-a^2\bigr)
\end{bmatrix}.
\]

Now set
\[
\mathbf T_n(\mu)\coloneqq
\bigl[\mathbf t(\zeta_1)\ \mathbf t(\zeta_2)\ \mathbf t(\zeta_3)\ \mathbf t(\zeta_4)\bigr]
\in\mathbb C^{4\times 4},
\]
\[
\mathbf N_n(\mu)\coloneqq
\bigl[\mathbf n(\zeta_1)\ \mathbf n(\zeta_2)\ \mathbf n(\zeta_3)\ \mathbf n(\zeta_4)\bigr]
\in\mathbb C^{4\times 4}.
\]

\begin{proposition}[Exact boundary Green matrix]\label{prop:four-root-boundary-green}
Assume \(\mu\notin\sigma(\mathbf H_{x,y})\), and assume \(p_\mu\) has four simple roots. Then
\[
\mathbf G_n(\mu)
\coloneqq
\mathbf P_\partial(\mathbf H_{x,y}-\mu\mathbf I_n)^{-1}\mathbf P_\partial^*
\]
is given exactly by
\[
\boxed{
\mathbf G_n(\mu)=\mathbf T_n(\mu)\,\mathbf N_n(\mu)^{-1}.
}
\]
If, in addition, \(\mathbf T_n(\mu)\) is invertible, then the boundary Weyl matrix
\[
\mathbf M_n(\mu)\coloneqq \mathbf N_n(\mu)\,\mathbf T_n(\mu)^{-1}
\]
is well defined and satisfies
\[
\boxed{
\mathbf M_n(\mu)=\mathbf G_n(\mu)^{-1}.
}
\]
\end{proposition}

\begin{proof}
Let
\[
\mathbf f\in\mathbb C^4,
\]
and solve
\[
(\mathbf H_{x,y}-\mu\mathbf I_n)\mathbf u=\mathbf P_\partial^* \mathbf f.
\]
Since the forcing is supported only at the boundary rows \(1,2,n-1,n\), the interior rows \(j=3,\dots,n-2\) satisfy the homogeneous recurrence from Proposition~\ref{prop:bulk-recurrence}. Therefore \(\mathbf u\) is a linear combination of the four root modes:
\[
\mathbf u=\sum_{\ell=1}^4 \gamma_\ell \mathbf u^{(\zeta_\ell)}
\]
for some
\[
\boldsymbol{\gamma}=
\begin{bmatrix}
\gamma_1\\
\gamma_2\\
\gamma_3\\
\gamma_4
\end{bmatrix}
\in\mathbb C^4.
\]

Applying \(\mathbf P_\partial\) gives
\[
\mathbf P_\partial\mathbf u=\mathbf T_n(\mu)\boldsymbol{\gamma},
\]
while applying \(\mathbf P_\partial(\mathbf H_{x,y}-\mu\mathbf I_n)\) gives
\[
\mathbf f=\mathbf N_n(\mu)\boldsymbol{\gamma}.
\]
Because \(\mu\notin\sigma(\mathbf H_{x,y})\), the boundary-value problem is uniquely solvable for every \(\mathbf f\), so \(\mathbf N_n(\mu)\) is invertible and
\[
\mathbf P_\partial\mathbf u=\mathbf T_n(\mu)\mathbf N_n(\mu)^{-1}\mathbf f.
\]
Since \(\mathbf f\) was arbitrary, this proves
\[
\mathbf G_n(\mu)=\mathbf T_n(\mu)\mathbf N_n(\mu)^{-1}.
\]

If \(\mathbf T_n(\mu)\) is invertible, then \(\mathbf G_n(\mu)\) is a product of two invertible matrices and therefore invertible, with
\[
\mathbf G_n(\mu)^{-1}
=
\mathbf N_n(\mu)\mathbf T_n(\mu)^{-1}
=
\mathbf M_n(\mu).
\]
\end{proof}

\begin{remark}
The four-root formula in Proposition~\ref{prop:four-root-boundary-green} is the exact finite-volume object. It remains valid all the way up to the finite-volume eigenvalue \(\lambda_1(x,y)\), provided the roots stay simple. This is the correct replacement for the purely subthreshold two-decaying-root ansatz.
\end{remark}

\subsection*{Subthreshold factorization through the two decaying roots}

Assume now that
\[
\mu<\mu_{\mathrm{sym}}(x,y),
\]
so that \(p_\mu\) has exactly two roots \(r_1,r_2\) in \(\mathbb D\), and the other two are
\[
\rho_1=\overline{r_1}^{\,-1},
\qquad
\rho_2=\overline{r_2}^{\,-1}.
\]
Set
\[
\sigma_1\coloneqq r_1+r_2,
\qquad
\sigma_2\coloneqq r_1r_2,
\qquad
\mathbf J\coloneqq
\begin{bmatrix}
0&1\\
1&0
\end{bmatrix}.
\]
Define
\[
F_k\coloneqq \frac{r_1^k-r_2^k}{r_1-r_2},
\qquad k\ge 0,
\]
with \(F_0=0\), \(F_1=1\), and
\[
\mathbf A_0\coloneqq \mathbf I_2,
\qquad
\mathbf A_k\coloneqq
\begin{bmatrix}
-\sigma_2 F_{k-1} & F_k\\
-\sigma_2 F_k & F_{k+1}
\end{bmatrix},
\qquad k\ge 1.
\]
Equivalently,
\[
\mathbf A_k=\mathbf C^k,
\qquad
\mathbf C\coloneqq
\begin{bmatrix}
0&1\\
-\sigma_2&\sigma_1
\end{bmatrix}.
\]
Also define
\[
\mathbf B_k\coloneqq \mathbf J\,\overline{\mathbf A_k}\,\mathbf J.
\]

\begin{lemma}\label{lem:block-propagators}
Let \(\mathbf u^L\) be a linear combination of the two left-decaying modes \(r_1^{j-1},r_2^{j-1}\). Then
\[
\begin{bmatrix}
u_{k+1}^L\\
u_{k+2}^L
\end{bmatrix}
=
\mathbf A_k
\begin{bmatrix}
u_1^L\\
u_2^L
\end{bmatrix}.
\]
Likewise, if \(\mathbf u^R\) is a linear combination of the two right-decaying modes
\(\rho_1^{j-n},\rho_2^{j-n}\), then
\[
\begin{bmatrix}
u_{n-k-1}^R\\
u_{n-k}^R
\end{bmatrix}
=
\mathbf B_k
\begin{bmatrix}
u_{n-1}^R\\
u_n^R
\end{bmatrix}.
\]
\end{lemma}

\begin{proof}
For \(\mathbf u^L\), the two-dimensional space spanned by \(r_1^{j-1},r_2^{j-1}\) satisfies the second-order recurrence
\[
u_{j+2}=\sigma_1 u_{j+1}-\sigma_2 u_j,
\]
hence the transfer matrix from \((u_1,u_2)\) to \((u_{k+1},u_{k+2})\) is \(\mathbf C^k\). The explicit formula for \(\mathbf A_k\) is the standard closed form for \(\mathbf C^k\). The right-decaying statement follows by conjugate reflection.
\end{proof}

For \(n\ge 5\), let
\[
\boldsymbol{\xi}\coloneqq
\begin{bmatrix}
u_1\\
u_2
\end{bmatrix},
\qquad
\boldsymbol{\eta}\coloneqq
\begin{bmatrix}
u_{n-1}\\
u_n
\end{bmatrix},
\qquad
\mathbf q_L\coloneqq
\begin{bmatrix}
u_3\\
u_4
\end{bmatrix},
\qquad
\mathbf q_R\coloneqq
\begin{bmatrix}
u_{n-3}\\
u_{n-2}
\end{bmatrix}.
\]
If \(\mathbf u=\mathbf u^L+\mathbf u^R\) is a homogeneous bulk solution, then
\[
\begin{bmatrix}\boldsymbol{\xi}\\ \boldsymbol{\eta}\end{bmatrix}
=
\begin{bmatrix}
\mathbf I_2 & \mathbf B_{n-2}\\
\mathbf A_{n-2} & \mathbf I_2
\end{bmatrix}
\begin{bmatrix}
\boldsymbol{\xi}^L\\
\boldsymbol{\eta}^R
\end{bmatrix},
\]
and
\[
\begin{bmatrix}
\mathbf q_L\\
\mathbf q_R
\end{bmatrix}
=
\begin{bmatrix}
\mathbf A_2 & \mathbf B_{n-4}\\
\mathbf A_{n-4} & \mathbf B_2
\end{bmatrix}
\begin{bmatrix}
\boldsymbol{\xi}^L\\
\boldsymbol{\eta}^R
\end{bmatrix}.
\]
Hence
\[
\begin{bmatrix}
\mathbf q_L\\
\mathbf q_R
\end{bmatrix}
=
\Gamma_n(\mu)
\begin{bmatrix}
\boldsymbol{\xi}\\
\boldsymbol{\eta}
\end{bmatrix},
\]
where
\[
\boxed{
\Gamma_n(\mu)\coloneqq
\begin{bmatrix}
\mathbf A_2 & \mathbf B_{n-4}\\
\mathbf A_{n-4} & \mathbf B_2
\end{bmatrix}
\begin{bmatrix}
\mathbf I_2 & \mathbf B_{n-2}\\
\mathbf A_{n-2} & \mathbf I_2
\end{bmatrix}^{-1}.
}
\]

Now set
\[
\mathbf L_0\coloneqq
\begin{bmatrix}
d_\mu-1 & -c\\
-\overline c & d_\mu
\end{bmatrix},
\qquad
\mathbf L_1\coloneqq
\begin{bmatrix}
a&0\\
-c&a
\end{bmatrix},
\]
\[
\mathbf R_0\coloneqq
\begin{bmatrix}
d_\mu & -c\\
-\overline c & d_\mu-a^2
\end{bmatrix},
\qquad
\mathbf R_1\coloneqq
\begin{bmatrix}
a&-\overline c\\
0&a
\end{bmatrix}.
\]

\begin{proposition}[Subthreshold \(2\times2\) block factorization]\label{prop:subthreshold-block-weyl}
Assume \(\mu<\mu_{\mathrm{sym}}(x,y)\), \(n\ge 5\), and that the matrix
\[
\begin{bmatrix}
\mathbf I_2 & \mathbf B_{n-2}\\
\mathbf A_{n-2} & \mathbf I_2
\end{bmatrix}
\]
is invertible. Then the boundary Weyl matrix \(\mathbf M_n(\mu)\) from Proposition~\ref{prop:four-root-boundary-green} is exactly
\[
\boxed{
\mathbf M_n(\mu)
=
\begin{bmatrix}
\mathbf L_0&0\\
0&\mathbf R_0
\end{bmatrix}
+
\begin{bmatrix}
\mathbf L_1&0\\
0&\mathbf R_1
\end{bmatrix}
\Gamma_n(\mu).
}
\]
\end{proposition}

\begin{proof}
For a homogeneous bulk solution \(\mathbf u\), the first two boundary rows are
\[
\begin{bmatrix}
f_1\\
f_2
\end{bmatrix}
=
\mathbf L_0\,\boldsymbol{\xi}+\mathbf L_1\,\mathbf q_L,
\]
and the last two are
\[
\begin{bmatrix}
f_{n-1}\\
f_n
\end{bmatrix}
=
\mathbf R_0\,\boldsymbol{\eta}+\mathbf R_1\,\mathbf q_R.
\]
Substitute
\[
\begin{bmatrix}
\mathbf q_L\\
\mathbf q_R
\end{bmatrix}
=
\Gamma_n(\mu)
\begin{bmatrix}
\boldsymbol{\xi}\\
\boldsymbol{\eta}
\end{bmatrix}.
\]
This gives exactly the displayed block formula.
\end{proof}

\begin{remark}
Proposition~\ref{prop:subthreshold-block-weyl} is the exact two-decaying-root factorization suggested by the numerical experiments. It is very useful below the symbol floor, but Proposition~\ref{prop:symbol-floor} shows that it is not, by itself, enough to reach \(\mu=\lambda_1(x,y)\) in finite volume. For that, one must use the four-root matrix \(\mathbf G_n(\mu)\), or the corresponding Weyl matrix \(\mathbf M_n(\mu)\) when the latter is defined.
\end{remark}

\subsection*{A boundary-Weyl derivative formula for \(\lambda_1(x,y)\)}

Assume now that \(\lambda_1(x,y)\) is simple, and let \(\boldsymbol{\psi}(x,y)\) be a normalized eigenvector of
\(\mathbf H_{x,y}\) for \(\lambda_1(x,y)\):
\[
\mathbf H_{x,y}\boldsymbol{\psi}=\lambda_1(x,y)\boldsymbol{\psi},
\qquad
\|\boldsymbol{\psi}\|_2=1.
\]
Set
\[
\mathbf b(x,y)\coloneqq \mathbf P_\partial \boldsymbol{\psi}(x,y)\in\mathbb C^4.
\]

\begin{proposition}[Boundary Weyl derivative identity]\label{prop:boundary-weyl-derivative}
Assume that \(p_{\lambda_1(x,y)}\) has four simple roots and that \(\mathbf T_n(\mu)\) is invertible for \(\mu\) in a neighborhood of \(\lambda_1(x,y)\). Let
\[
\mathbf M_n(\mu)\coloneqq \mathbf N_n(\mu)\mathbf T_n(\mu)^{-1}.
\]
Then:
\begin{enumerate}
\item
\[
\mathbf M_n(\lambda_1(x,y))\,\mathbf b(x,y)=0.
\]
\item
\[
\mathbf b(x,y)^*\,\partial_\mu \mathbf M_n(\lambda_1(x,y))\,\mathbf b(x,y)=-1.
\]
\item Consequently,
\[
\boxed{
\partial_y \lambda_1(x,y)
=
\mathbf b(x,y)^*\,\partial_y \mathbf M_n(\lambda_1(x,y))\,\mathbf b(x,y).
}
\]
\end{enumerate}
Thus the vertical monotonicity problem for \(s(x,y)^2=\lambda_1(x,y)\) is exactly the positivity problem
\[
\mathbf b^*\,\partial_y \mathbf M_n(\lambda_1(x,y),y)\,\mathbf b\ge 0
\qquad
\text{on }\ker \mathbf M_n(\lambda_1(x,y),y).
\]
\end{proposition}

\begin{proof}
Let
\[
\mathbf G_n(\mu)=\mathbf P_\partial(\mathbf H_{x,y}-\mu\mathbf I_n)^{-1}\mathbf P_\partial^*.
\]
By the spectral theorem,
\[
\mathbf G_n(\mu)
=
\sum_{k=1}^n
\frac{\mathbf b_k\mathbf b_k^*}{\lambda_k-\mu},
\qquad
\mathbf b_k\coloneqq \mathbf P_\partial \boldsymbol{\psi}_k,
\]
where \(\{\boldsymbol{\psi}_k\}\) is an orthonormal eigenbasis of \(\mathbf H_{x,y}\) and \(\lambda_1=\lambda_1(x,y)\). Since \(\lambda_1\) is simple and \(\mathbf b=\mathbf P_\partial \boldsymbol{\psi}\),
\[
\mathbf G_n(\mu)
=
\frac{\mathbf b\,\mathbf b^*}{\lambda_1-\mu}
+
\mathbf G_n^{\mathrm{reg}}(\mu)
\]
near \(\mu=\lambda_1\).

By Proposition~\ref{prop:four-root-boundary-green}, \(\mathbf G_n(\mu)=\mathbf T_n(\mu)\mathbf N_n(\mu)^{-1}\), and by the invertibility assumption on \(\mathbf T_n(\mu)\), the matrix
\[
\mathbf M_n(\mu)=\mathbf N_n(\mu)\mathbf T_n(\mu)^{-1}
\]
is defined near \(\lambda_1\) and satisfies
\[
\mathbf M_n(\mu)\mathbf G_n(\mu)=\mathbf I_4
\]
for \(\mu\neq \lambda_1\) close to \(\lambda_1\). Comparing the coefficient of \((\lambda_1-\mu)^{-1}\) gives
\[
\mathbf M_n(\lambda_1)\mathbf b=0.
\]
Using instead \(\mathbf G_n(\mu)\mathbf M_n(\mu)=\mathbf I_4\) and applying both sides to \(\mathbf b\), we obtain
\[
\left(\frac{\mathbf b\,\mathbf b^*}{\lambda_1-\mu}+\mathbf G_n^{\mathrm{reg}}(\mu)\right)\mathbf M_n(\mu)\mathbf b=\mathbf b.
\]
Since \(\mathbf M_n(\lambda_1)\mathbf b=0\), we have
\[
\mathbf M_n(\mu)\mathbf b
=
(\mu-\lambda_1)\,\partial_\mu \mathbf M_n(\lambda_1)\mathbf b+o(\mu-\lambda_1),
\]
and therefore
\[
-\mathbf b\,\mathbf b^*\,\partial_\mu \mathbf M_n(\lambda_1)\mathbf b=\mathbf b.
\]
Left-multiplying by \(\mathbf b^*\) yields
\[
-\|\mathbf b\|_2^2\,\mathbf b^*\,\partial_\mu \mathbf M_n(\lambda_1)\mathbf b
=
\|\mathbf b\|_2^2,
\]
hence
\[
\mathbf b^*\,\partial_\mu \mathbf M_n(\lambda_1)\mathbf b=-1.
\]

Now differentiate the kernel relation
\[
\mathbf M_n(\lambda_1(x,y),y)\,\mathbf b(x,y)=0
\]
with respect to \(y\):
\[
\Bigl(\partial_\mu \mathbf M_n(\lambda_1,y)\,\partial_y\lambda_1
+\partial_y\mathbf M_n(\lambda_1,y)\Bigr)\mathbf b
+
\mathbf M_n(\lambda_1,y)\,\partial_y\mathbf b
=0.
\]
Left-multiplying by \(\mathbf b^*\), and using \(\mathbf M_n(\lambda_1)\mathbf b=0\) and Hermiticity of \(\mathbf M_n(\lambda_1)\), yields
\[
\mathbf b^*\,\partial_\mu \mathbf M_n(\lambda_1)\,\mathbf b\;\partial_y\lambda_1
+
\mathbf b^*\,\partial_y\mathbf M_n(\lambda_1)\,\mathbf b
=0.
\]
Using \(\mathbf b^*\,\partial_\mu \mathbf M_n(\lambda_1)\,\mathbf b=-1\), we obtain
\[
\partial_y\lambda_1
=
\mathbf b^*\,\partial_y\mathbf M_n(\lambda_1)\,\mathbf b.
\]
\end{proof}

\begin{corollary}[Root-coordinate form]\label{cor:root-coordinate-derivative}
Assume the hypotheses of Proposition~\ref{prop:boundary-weyl-derivative}. Let
\[
\boldsymbol{\gamma}(x,y)\in\ker \mathbf N_n(\lambda_1(x,y))
\]
be chosen so that
\[
\mathbf b(x,y)=\mathbf T_n(\lambda_1(x,y))\,\boldsymbol{\gamma}(x,y).
\]
Then
\[
\boxed{
\partial_y\lambda_1(x,y)
=
\boldsymbol{\gamma}(x,y)^*\,\mathbf T_n(\lambda_1)^*
\bigl(\partial_y \mathbf N_n(\lambda_1)\bigr)\,
\boldsymbol{\gamma}(x,y).
}
\]
Moreover,
\[
\boxed{
\boldsymbol{\gamma}(x,y)^*\,\mathbf T_n(\lambda_1)^*
\bigl(\partial_\mu \mathbf N_n(\lambda_1)\bigr)\,
\boldsymbol{\gamma}(x,y)
=
-1.
}
\]
\end{corollary}

\begin{proof}
Since
\[
\mathbf M_n(\mu)=\mathbf N_n(\mu)\mathbf T_n(\mu)^{-1},
\]
and \(\mathbf b=\mathbf T_n(\lambda_1)\boldsymbol{\gamma}\) with \(\mathbf N_n(\lambda_1)\boldsymbol{\gamma}=0\), we have
\[
\mathbf M_n(\lambda_1)\mathbf b=0,
\]
and
\[
\partial_y\mathbf M_n(\lambda_1)\,\mathbf b
=
\bigl(\partial_y\mathbf N_n(\lambda_1)\bigr)\boldsymbol{\gamma},
\]
because the term involving \(\partial_y(\mathbf T_n^{-1})\) is annihilated by
\(\mathbf N_n(\lambda_1)\boldsymbol{\gamma}\). Therefore
\[
\mathbf b^*\,\partial_y\mathbf M_n(\lambda_1)\,\mathbf b
=
\boldsymbol{\gamma}^*\,\mathbf T_n(\lambda_1)^*\,
\bigl(\partial_y\mathbf N_n(\lambda_1)\bigr)\boldsymbol{\gamma}.
\]
Now apply Proposition~\ref{prop:boundary-weyl-derivative}. The \(\mu\)-derivative formula is identical.
\end{proof}

\subsection*{Back to the reduced-resolvent criterion}

Let \(\widehat{\mathbf C}_x\) be the \(4\times4\) boundary matrix corresponding to the operator \(\mathbf C_x\) from Proposition~\ref{prop:B-commutator}:
\[
\widehat{\mathbf C}_x=
\begin{bmatrix}
-2x&1&0&0\\
1&0&0&0\\
0&0&0&-a\\
0&0&-a&2x
\end{bmatrix},
\qquad
\mathbf C_x=\mathbf P_\partial^*\,\widehat{\mathbf C}_x\,\mathbf P_\partial.
\]
Set
\[
\mathbf q(x,y)\coloneqq \widehat{\mathbf C}_x\,\mathbf b(x,y)\in\mathbb C^4.
\]
Because \(\widehat{\mathbf C}_x\) is Hermitian and
\[
[\mathbf H_{x,y},\mathbf L]=i(1+a)\mathbf P_\partial^* \widehat{\mathbf C}_x \mathbf P_\partial
\]
by Proposition~\ref{prop:B-commutator}, we have
\[
\mathbf b^*\mathbf q
=
\boldsymbol{\psi}^*\,\mathbf P_\partial^* \widehat{\mathbf C}_x \mathbf P_\partial\,\boldsymbol{\psi}
=
\frac{1}{i(1+a)}\boldsymbol{\psi}^*[\mathbf H_{x,y},\mathbf L]\boldsymbol{\psi}
=0.
\]
Hence the pole of the boundary resolvent cancels on the vector \(\mathbf q\).

\begin{proposition}[Exact \(4\times4\) scalar Stieltjes reduction]\label{prop:4x4-stieltjes}
Define
\[
g_{x,y}(\mu)\coloneqq
\mathbf q(x,y)^*\,\mathbf G_n(\mu)\,\mathbf q(x,y)
=
\mathbf q(x,y)^*\,\mathbf T_n(\mu)\mathbf N_n(\mu)^{-1}\mathbf q(x,y),
\qquad
\mu\notin\sigma(\mathbf H_{x,y}).
\]
Then \(g_{x,y}(\mu)\) extends holomorphically through \(\mu=\lambda_1(x,y)\), and
\[
\boxed{
\left\langle
\mathbf C_x\boldsymbol{\psi}(x,y),\,
(\mathbf H_{x,y}-\lambda_1(x,y)\mathbf I_n)^{-3}_{\perp}\,
\mathbf C_x\boldsymbol{\psi}(x,y)
\right\rangle
=
\frac12\,g_{x,y}''\bigl(\lambda_1(x,y)\bigr).
}
\]
Consequently, the reduced-resolvent convexity inequality from Proposition~\ref{prop:boundary-green} is equivalent to
\[
\boxed{
(1-a^2)^2\,g_{x,y}''\bigl(\lambda_1(x,y)\bigr)\le 2.
}
\]
\end{proposition}

\begin{proof}
Since \(\mathbf q\perp \mathbf b\), the pole term
\[
\frac{\mathbf b\,\mathbf b^*}{\lambda_1-\mu}
\]
in \(\mathbf G_n(\mu)\) does not contribute to \(\mathbf q^*\mathbf G_n(\mu)\mathbf q\). Thus \(g_{x,y}\) is regular at \(\mu=\lambda_1\).

Let \(\{\boldsymbol{\psi}_k\}_{k=1}^n\) be an orthonormal eigenbasis of \(\mathbf H_{x,y}\), with
\(\boldsymbol{\psi}_1=\boldsymbol{\psi}\) and \(\lambda_1=\lambda_1(x,y)\). Set
\[
\mathbf b_k\coloneqq \mathbf P_\partial \boldsymbol{\psi}_k\in\mathbb C^4.
\]
Then
\[
\mathbf G_n(\mu)
=
\sum_{k=1}^n \frac{\mathbf b_k\mathbf b_k^*}{\lambda_k-\mu},
\]
and therefore
\[
g_{x,y}(\mu)
=
\sum_{k=1}^n \frac{|\mathbf q^*\mathbf b_k|^2}{\lambda_k-\mu}.
\]
Because \(\mathbf q^*\mathbf b_1=\mathbf q^*\mathbf b=0\), we get
\[
g_{x,y}(\mu)
=
\sum_{k=2}^n \frac{|\mathbf q^*\mathbf b_k|^2}{\lambda_k-\mu},
\]
and hence
\[
\frac12\,g_{x,y}''(\lambda_1)
=
\sum_{k=2}^n
\frac{|\mathbf q^*\mathbf b_k|^2}{(\lambda_k-\lambda_1)^3}.
\]
Now
\[
\mathbf q^*\mathbf b_k
=
\mathbf b^*\widehat{\mathbf C}_x \mathbf b_k
=
\boldsymbol{\psi}^* \mathbf P_\partial^* \widehat{\mathbf C}_x \mathbf P_\partial \boldsymbol{\psi}_k
=
\langle \mathbf C_x\boldsymbol{\psi},\boldsymbol{\psi}_k\rangle,
\]
so the previous sum is exactly
\[
\left\langle
\mathbf C_x\boldsymbol{\psi},
(\mathbf H_{x,y}-\lambda_1\mathbf I_n)^{-3}_{\perp}
\mathbf C_x\boldsymbol{\psi}
\right\rangle.
\]
The last equivalence is Proposition~\ref{prop:boundary-green}.
\end{proof}

\begin{remark}[Where the proof now stands]
The route through the bulk current has now been replaced by an exact \(4\times4\) boundary-Weyl problem.

\smallskip

\noindent
The vertical monotonicity of \(\lambda_1(x,y)=s(x,y)^2\) is equivalent to
\[
\mathbf b^*\,\partial_y\mathbf M_n(\lambda_1(x,y),y)\,\mathbf b\ge 0
\qquad
\text{for }\mathbf b\in\ker \mathbf M_n(\lambda_1(x,y),y),
\]
whenever the matrix \(\mathbf M_n\) is defined near \(\lambda_1(x,y)\). Equivalently, in root coordinates,
\[
\boldsymbol{\gamma}^*\,\mathbf T_n(\lambda_1)^*
\bigl(\partial_y\mathbf N_n(\lambda_1)\bigr)\boldsymbol{\gamma}\ge 0
\qquad
\text{for }\boldsymbol{\gamma}\in\ker \mathbf N_n(\lambda_1).
\]
And the reduced-resolvent bound is equivalent to the scalar inequality
\[
(1-a^2)^2\,g_{x,y}''(\lambda_1(x,y))\le 2
\]
for the explicit \(4\times4\) Stieltjes function
\[
g_{x,y}(\mu)=\mathbf q^*\,\mathbf T_n(\mu)\mathbf N_n(\mu)^{-1}\mathbf q.
\]

\smallskip

\noindent
So the analytic heart of the problem is now compressed to an explicit \(4\times4\) Evans/Weyl matrix, valid up to the actual finite-volume eigenvalue. The remaining missing step is a sign estimate for \(\partial_y\mathbf M_n\) on its kernel, or equivalently the second-derivative bound above for \(g_{x,y}\).
\end{remark}

\section{Further progress toward Conjecture~\ref{conj:vertical-monotonicity}}

In this section we record several rigorous partial results for the monotonicity problem. The strongest one is a complete proof on the imaginary axis \(x=0\). We then prove a continuation result in the \(x\)-variable along simple branches on which the right singular vector has no zero entries.

For \(x\in\mathbb R\) and \(y>0\), write
\[
\mathbf B_{x,y}\coloneqq (x+iy)\mathbf I_n-\mathbf A_n(a),
\qquad
\mathbf H_x(y)\coloneqq \mathbf B_{x,y}^*\,\mathbf B_{x,y}.
\]
Whenever \(\lambda_1(x,y)=s(x,y)^2\) is simple, we choose a normalized right singular vector
\[
\mathbf v(x,y)=
\begin{bmatrix}
v_1(x,y)\\
\vdots\\
v_n(x,y)
\end{bmatrix}
\]
and a normalized left singular vector
\[
\mathbf u(x,y)=
\begin{bmatrix}
u_1(x,y)\\
\vdots\\
u_n(x,y)
\end{bmatrix}
\]
in the canonical gauge
\[
\mathbf B_{x,y}\mathbf v=s(x,y)\,\mathbf u,
\qquad
\mathbf B_{x,y}^*\mathbf u=s(x,y)\,\mathbf v,
\qquad
\mathbf u=\mathbf R\,\overline{\mathbf v}.
\]
Then the singular-vector equation reads componentwise
\[
(x+iy)v_j-a v_{j-1}-v_{j+1}=s(x,y)\,\overline{v_{n+1-j}},
\qquad
j=1,\dots,n,
\]
with \(v_0=v_{n+1}=0\).

\begin{lemma}[A real telescoping invariant]\label{lem:W-real}
Assume \(\lambda_1(x,y)=s(x,y)^2\) is simple and let \(\mathbf v\) be in the canonical gauge above. Define
\[
W_j\coloneqq v_{j+1}v_{n+1-j}-a\,v_jv_{n-j},
\qquad
j=0,\dots,n,
\]
with \(v_0=v_{n+1}=0\). Then
\[
W_0=0,
\qquad
W_j-W_{j-1}=s(x,y)\bigl(|v_j|^2-|v_{n+1-j}|^2\bigr)\in\mathbb R.
\]
In particular,
\[
W_j\in\mathbb R
\qquad
(j=0,\dots,n).
\]
\end{lemma}

\begin{proof}
Using the singular-vector equation at indices \(j\) and \(n+1-j\), one gets
\[
v_{j+1}+a v_{j-1}=(x+iy)v_j-s\,\overline{v_{n+1-j}},
\]
and
\[
v_{n+2-j}+a v_{n-j}=(x+iy)v_{n+1-j}-s\,\overline{v_j}.
\]
Subtracting,
\begin{align*}
W_j-W_{j-1}
&=
v_{n+1-j}\bigl(v_{j+1}+a v_{j-1}\bigr)
-
v_j\bigl(v_{n+2-j}+a v_{n-j}\bigr) \\
&=
v_{n+1-j}\bigl((x+iy)v_j-s\,\overline{v_{n+1-j}}\bigr)
-
v_j\bigl((x+iy)v_{n+1-j}-s\,\overline{v_j}\bigr) \\
&=
s\bigl(|v_j|^2-|v_{n+1-j}|^2\bigr)\in\mathbb R.
\end{align*}
Since \(W_0=0\), it follows by induction that \(W_j\in\mathbb R\) for all \(j\).
\end{proof}

\begin{proposition}[Imaginary-axis structure and monotonicity]\label{prop:x=0-monotonicity}
For every \(y>0\), the eigenvalue \(\lambda_1(0,y)=s(0,y)^2\) is simple. Moreover, the corresponding canonical right singular vector \(\mathbf v(0,y)\) can be chosen so that
\[
\Im\!\left(\frac{v_{j+1}(0,y)}{v_j(0,y)}\right)>0,
\qquad
j=1,\dots,n-1,
\]
and
\[
\Im\!\left(\frac{\overline{v_{n+1-j}(0,y)}}{v_j(0,y)}\right)>0,
\qquad
j=1,\dots,n.
\]
Consequently,
\[
\frac{\partial}{\partial y}\bigl(s(0,y)^2\bigr)\ge 0
\qquad (y>0),
\]
and hence Conjecture~\ref{conj:vertical-monotonicity} holds for \(x=0\).
\end{proposition}

\begin{proof}
Set
\[
\mathbf B_{0,y}=iy\,\mathbf I_n-\mathbf A_n(a),
\qquad
\mathbf H_0(y)=\mathbf B_{0,y}^*\,\mathbf B_{0,y}.
\]
Let
\[
\mathbf D\coloneqq \operatorname{diag}\bigl(1,i,i^2,\dots,i^{n-1}\bigr),
\qquad
\widehat{\mathbf H}_y\coloneqq \mathbf D^*\,\mathbf H_0(y)\,\mathbf D.
\]
A direct computation shows that \(\widehat{\mathbf H}_y\) is real symmetric, with diagonal entries
\[
(\widehat{\mathbf H}_y)_{11}=a^2+y^2,
\qquad
(\widehat{\mathbf H}_y)_{jj}=1+a^2+y^2\quad (2\le j\le n-1),
\qquad
(\widehat{\mathbf H}_y)_{nn}=1+y^2,
\]
and nonzero off-diagonal entries
\[
(\widehat{\mathbf H}_y)_{j,j+1}=-(1-a)y<0,
\qquad
(\widehat{\mathbf H}_y)_{j,j+2}=-a<0,
\]
whenever the indices are defined. Thus \(\widehat{\mathbf H}_y\) is a real symmetric irreducible \(5\)-diagonal \(Z\)-matrix.

Choose \(\Lambda>\max_{1\le j\le n}(\widehat{\mathbf H}_y)_{jj}\). Then \(\Lambda\mathbf I_n-\widehat{\mathbf H}_y\) is entrywise nonnegative and irreducible, so by Perron--Frobenius its spectral radius is simple and has a strictly positive eigenvector. Equivalently, the smallest eigenvalue of \(\widehat{\mathbf H}_y\) is simple and has an eigenvector
\[
\mathbf w=
\begin{bmatrix}
w_1\\
\vdots\\
w_n
\end{bmatrix},
\qquad
w_j>0\quad (j=1,\dots,n).
\]
Since \(\mathbf H_0(y)\) and \(\widehat{\mathbf H}_y\) are unitarily similar, the normalized right singular vector of \(\mathbf B_{0,y}\) for \(s(0,y)\) may be written as
\[
v_j=\beta\,i^{\,j-1}w_j,
\qquad
j=1,\dots,n,
\]
for some unimodular \(\beta\).

Now insert this into the singular-vector equation
\[
iy\,v_j-a v_{j-1}-v_{j+1}=s(0,y)\,\overline{v_{n+1-j}}.
\]
After dividing by \(\beta i^j\), one obtains
\[
w_{j+1}-y w_j-a w_{j-1}=-\gamma\,s(0,y)\,w_{n+1-j},
\qquad
\gamma\coloneqq \frac{\overline\beta}{\beta}(-i)^n.
\]
At \(j=n\), this gives
\[
y w_n+a w_{n-1}=\gamma\,s(0,y)\,w_1.
\]
The left-hand side is strictly positive, and \(|\gamma|=1\), so necessarily \(\gamma=1\).

With this normalization,
\[
r_j\coloneqq \frac{v_{j+1}}{v_j}=i\,\frac{w_{j+1}}{w_j},
\qquad
j=1,\dots,n-1,
\]
and hence
\[
\Im r_j>0,
\qquad
j=1,\dots,n-1.
\]
Also,
\[
p_j\coloneqq \frac{\overline{v_{n+1-j}}}{v_j}=i\,\frac{w_{n+1-j}}{w_j},
\qquad
j=1,\dots,n,
\]
so
\[
\Im p_j>0
\qquad
(j=1,\dots,n).
\]

By Proposition~\ref{prop:current-identity},
\[
J_j-aJ_{j-1}=y|v_j|^2+s(0,y)\,\Im(v_jv_{n+1-j}).
\]
But
\[
\Im p_j
=
-\frac{\Im(v_jv_{n+1-j})}{|v_j|^2},
\]
so \(\Im p_j>0\) implies
\[
\Im(v_jv_{n+1-j})<0.
\]
Hence
\[
J_j-aJ_{j-1}\le y|v_j|^2.
\]
Inductively,
\[
J_j\le y\sum_{k=1}^j a^{\,j-k}|v_k|^2,
\]
and therefore
\[
(1-a)\sum_{j=1}^{n-1}J_j
\le
y\sum_{k=1}^{n-1}(1-a^{\,n-k})|v_k|^2
\le
y.
\]
Using Proposition~\ref{prop:current-identity} again,
\[
\frac{\partial}{\partial y}\bigl(s(0,y)^2\bigr)
=
2y-2(1-a)\sum_{j=1}^{n-1}J_j
\ge 0.
\]
\end{proof}

\begin{lemma}[Boundary entries never vanish]\label{lem:boundary-entries-nonzero}
Assume \(\lambda_1(x,y)=s(x,y)^2\) is simple, \(y>0\), and let \(\mathbf v\) be in the canonical gauge. Then
\[
v_1\neq 0,
\qquad
v_n\neq 0.
\]
\end{lemma}

\begin{proof}
We first prove \(v_n\neq 0\). Assume \(v_n=0\). The case \(n=2\) is immediate, so assume \(n\ge 3\).

If also \(v_1=0\), then the equations at \(j=1\) and \(j=n\) imply \(v_2=0\) and \(v_{n-1}=0\). Repeatedly using the singular-vector equations then yields \(v_j=0\) for every \(j\), contradicting \(\|\mathbf v\|_2=1\). Hence \(v_1\neq 0\).

From the singular-vector equation at \(j=1\),
\[
(x+iy)v_1-v_2=s\,\overline{v_n}=0,
\]
hence
\[
v_2=(x+iy)v_1.
\]
At \(j=n-1\),
\[
(x+iy)v_{n-1}-a v_{n-2}-v_n=s\,\overline{v_2},
\]
so
\[
(x+iy)v_{n-1}-a v_{n-2}=s\,\overline{x+iy}\,\overline{v_1}.
\]
Now
\[
W_{n-2}=v_{n-1}v_2-a v_{n-2}v_1
=
v_1\bigl((x+iy)v_{n-1}-a v_{n-2}\bigr)
=
s\,\overline{x+iy}\,|v_1|^2.
\]
By Lemma~\ref{lem:W-real}, \(W_{n-2}\in\mathbb R\), but
\[
\Im\!\bigl(s\,\overline{x+iy}\,|v_1|^2\bigr)=-s\,y\,|v_1|^2\neq 0,
\]
a contradiction. Thus \(v_n\neq 0\).

Now assume \(v_1=0\). By the first part, \(v_n\neq 0\). From \(j=1\),
\[
-v_2=s\,\overline{v_n},
\qquad\text{so}\qquad
v_2=-s\,\overline{v_n}.
\]
From \(j=n\),
\[
(x+iy)v_n-a v_{n-1}=s\,\overline{v_1}=0,
\qquad\text{so}\qquad
a v_{n-1}=(x+iy)v_n.
\]
At \(j=2\),
\[
(x+iy)v_2-v_3=s\,\overline{v_{n-1}}.
\]
Hence
\[
v_3=(x+iy)v_2-s\,\overline{v_{n-1}}.
\]
Therefore
\[
W_2=v_3v_n-a v_2v_{n-1}
=
-s\,\overline{v_{n-1}}\,v_n
=
-\frac{s}{a}\,\overline{x+iy}\,|v_n|^2.
\]
Again \(W_2\in\mathbb R\) by Lemma~\ref{lem:W-real}, while
\[
\Im\!\left(-\frac{s}{a}\,\overline{x+iy}\,|v_n|^2\right)
=
\frac{s\,y}{a}|v_n|^2\neq 0.
\]
This contradiction proves \(v_1\neq 0\).
\end{proof}

\begin{proposition}[Persistence of \(\Im r_j>0\) along nonvanishing simple branches]\label{prop:r-positivity-branch}
Fix \(y>0\), and let \(I\subset\mathbb R\) be an open interval on which \(\lambda_1(x,y)=s(x,y)^2\) is simple. Let
\[
\mathbf v(x)=
\begin{bmatrix}
v_1(x)\\
\vdots\\
v_n(x)
\end{bmatrix}
\]
be the corresponding analytic right singular vector in the canonical gauge. Let \(J\) be a connected component of
\[
\{x\in I:\ v_1(x)\cdots v_n(x)\neq 0\}.
\]
Assume that for some \(x_0\in J\),
\[
\Im\!\left(\frac{v_{j+1}(x_0)}{v_j(x_0)}\right)>0,
\qquad
j=1,\dots,n-1.
\]
Then for every \(x\in J\),
\[
\Im\!\left(\frac{v_{j+1}(x)}{v_j(x)}\right)>0,
\qquad
j=1,\dots,n-1.
\]
Equivalently, if we define
\[
r_j(x)\coloneqq \frac{v_{j+1}(x)}{v_j(x)},
\qquad
\rho_j(x)\coloneqq r_{n-j}(x),
\]
then
\[
\Im r_j(x)>0,
\qquad
\Im \rho_j(x)>0,
\qquad
j=1,\dots,n-1,
\quad x\in J.
\]
\end{proposition}

\begin{proof}
Since \(v_j(x)\neq 0\) on \(J\), we may write
\[
v_j(x)=t_j(x)e^{i\phi_j(x)},
\qquad
t_j(x)>0,
\]
and define
\[
\delta_j(x)\coloneqq \phi_{j+1}(x)-\phi_j(x),
\qquad
j=1,\dots,n-1.
\]
Then
\[
r_j(x)=\frac{t_{j+1}(x)}{t_j(x)}e^{i\delta_j(x)},
\]
so \(\Im r_j(x)>0\) is equivalent to
\[
\delta_j(x)\in(0,\pi).
\]

Set
\[
c\coloneqq (1+a)x-i(1-a)y.
\]
Since \(y>0\), we may write
\[
c=|c|e^{-i\theta},
\qquad
\theta\in(0,\pi).
\]
A direct expansion of the quadratic form of \(\mathbf H_x(y)\) shows that
\[
\langle \mathbf v,\mathbf H_x(y)\mathbf v\rangle
=
\sum_{j=1}^n d_j\,t_j^2
+
\mathcal P_x(\delta_1,\dots,\delta_{n-1}),
\]
where
\[
d_1=a^2+x^2+y^2,
\qquad
d_j=1+a^2+x^2+y^2\quad(2\le j\le n-1),
\qquad
d_n=1+x^2+y^2,
\]
and
\[
\mathcal P_x(\delta_1,\dots,\delta_{n-1})
=
-2|c|\sum_{j=1}^{n-1} t_j t_{j+1}\cos(\delta_j-\theta)
+
2a\sum_{j=1}^{n-2} t_j t_{j+2}\cos(\delta_j+\delta_{j+1}).
\]
For fixed amplitudes \(t_1,\dots,t_n\), the phase vector
\[
\bigl[\delta_1(x),\dots,\delta_{n-1}(x)\bigr]
\]
must minimize \(\mathcal P_x\), because \(\mathbf v(x)\) is a ground state of \(\mathbf H_x(y)\).

Let
\[
E\coloneqq \{x\in J:\ \delta_j(x)\in(0,\pi)\ \text{for all }j=1,\dots,n-1\}.
\]
By hypothesis, \(x_0\in E\), so \(E\neq\varnothing\). Clearly \(E\) is open in \(J\). We show that \(E\) is also closed in \(J\). Let \(x_*\in \overline E\cap J\). Then
\[
\delta_j(x_*)\in[0,\pi]
\qquad
(j=1,\dots,n-1).
\]
If every \(\delta_j(x_*)\in(0,\pi)\), then \(x_*\in E\). So suppose some \(\delta_k(x_*)\in\{0,\pi\}\).

Consider the tail-rotation variation
\[
\phi_m\mapsto \phi_m+\varepsilon,
\qquad
m\ge k+1.
\]
This changes only \(\delta_k\), replacing it by \(\delta_k+\varepsilon\). Therefore
\[
f_k(\varepsilon)\coloneqq
\mathcal P_x(\delta_1,\dots,\delta_{k-1},\delta_k+\varepsilon,\delta_{k+1},\dots,\delta_{n-1})
\]
satisfies \(f_k'(0)=0\) at the minimizing phase vector. A direct differentiation gives
\begin{align*}
f_k'(0)
&=
2|c|\,t_k t_{k+1}\sin(\delta_k-\theta) \\
&\quad
-2a\,\mathbf 1_{\{k\ge 2\}}\,t_{k-1}t_{k+1}\sin(\delta_{k-1}+\delta_k) \\
&\quad
-2a\,\mathbf 1_{\{k\le n-2\}}\,t_k t_{k+2}\sin(\delta_k+\delta_{k+1}).
\end{align*}

If \(\delta_k(x_*)=0\), then \(\theta\in(0,\pi)\) implies
\[
\sin(\delta_k-\theta)=\sin(-\theta)<0,
\]
while the other two terms are \(\ge 0\), because \(\delta_{k-1}(x_*),\delta_{k+1}(x_*)\in[0,\pi]\). Hence \(f_k'(0)<0\), contradicting minimality.

If \(\delta_k(x_*)=\pi\), then
\[
\sin(\delta_k-\theta)=\sin(\pi-\theta)>0,
\]
and again the other two terms are \(\ge 0\). Hence \(f_k'(0)>0\), again a contradiction.

Therefore no \(\delta_k\) can hit \(0\) or \(\pi\), so \(x_*\in E\). Thus \(E\) is closed in \(J\). Since \(J\) is connected and \(E\subset J\) is nonempty, open, and closed, we conclude \(E=J\). This proves \(\Im r_j(x)>0\) for all \(x\in J\), and the statement for \(\rho_j(x)=r_{n-j}(x)\) follows immediately.
\end{proof}

\begin{lemma}[Exact angular identity]\label{lem:angular-identity}
Assume \(\lambda_1(x,y)=s(x,y)^2\) is simple and \(\mathbf v\) is in the canonical gauge. Let
\[
r_j\coloneqq \frac{v_{j+1}}{v_j},
\qquad
\rho_j\coloneqq r_{n-j},
\qquad
p_j\coloneqq \frac{\overline{v_{n+1-j}}}{v_j},
\]
and define
\[
\alpha_j\coloneqq (x+iy)-\frac{a}{r_{j-1}},
\qquad
\beta_j\coloneqq (x-iy)-\overline{\rho_{j-1}},
\qquad
j=1,\dots,n-1,
\]
with the convention \(r_0=\infty\). Then
\[
\Im\!\bigl(\alpha_j\overline{p_j}\bigr)
=
-\Im\!\bigl(\beta_j p_j\bigr),
\qquad
j=1,\dots,n-1.
\]
\end{lemma}

\begin{proof}
From the singular-vector equation at index \(j\),
\[
\biggl((x+iy)-\frac{a}{r_{j-1}}\biggr)v_j
=
v_{j+1}+s\,\overline{v_{n+1-j}},
\]
so
\[
\alpha_j\overline{p_j}
=
\frac{v_{j+1}v_{n+1-j}+s\,|v_{n+1-j}|^2}{|v_j|^2}.
\]
Hence
\[
\Im\!\bigl(\alpha_j\overline{p_j}\bigr)
=
\frac{\Im\!\bigl(v_{j+1}v_{n+1-j}\bigr)}{|v_j|^2}.
\]

Now conjugate the singular-vector equation at index \(n+1-j\):
\[
(x-iy)\,\overline{v_{n+1-j}}-a\,\overline{v_{n-j}}-\overline{v_{n+2-j}}
=
s\,v_j.
\]
Therefore
\[
\beta_j p_j
=
\bigl((x-iy)-\overline{\rho_{j-1}}\bigr)\frac{\overline{v_{n+1-j}}}{v_j}
=
s+a\,\frac{\overline{v_{n-j}}}{v_j},
\]
and thus
\[
\Im(\beta_j p_j)
=
-\frac{a\,\Im(v_jv_{n-j})}{|v_j|^2}.
\]

Finally, Lemma~\ref{lem:W-real} gives
\[
v_{j+1}v_{n+1-j}-a\,v_jv_{n-j}\in\mathbb R,
\]
so
\[
\Im\!\bigl(v_{j+1}v_{n+1-j}\bigr)
=
a\,\Im(v_jv_{n-j}).
\]
Substituting into the two displays above yields
\[
\Im\!\bigl(\alpha_j\overline{p_j}\bigr)
=
-\Im\!\bigl(\beta_j p_j\bigr).
\]
\end{proof}

\begin{proposition}[Persistence of \(\Im p_j>0\) on nonvanishing simple branches]\label{prop:p-positivity-branch}
Fix \(y>0\), and let \(J\) be as in Proposition~\ref{prop:r-positivity-branch}. Assume that on \(J\),
\[
\Im r_j(x)>0,
\qquad
\Im \rho_j(x)>0,
\qquad
j=1,\dots,n-1,
\]
and that for some \(x_0\in J\),
\[
\Im p_j(x_0)>0,
\qquad
j=1,\dots,n.
\]
Then
\[
\Im p_j(x)>0,
\qquad
j=1,\dots,n,
\quad x\in J.
\]
\end{proposition}

\begin{proof}
Fix \(j\in\{1,\dots,n\}\), and suppose that \(\Im p_j\) vanishes somewhere on \(J\). Since \(p_j\) depends continuously on \(x\), there exists \(x_*\in J\) with
\[
p_j(x_*)\in\mathbb R.
\]

We first note that \(p_j(x_*)\neq 0\). Indeed, by definition,
\[
p_j=\frac{\overline{v_{n+1-j}}}{v_j},
\]
and \(v_j,v_{n+1-j}\neq 0\) on \(J\) by construction.

At \(x_*\), the number \(p_j(x_*)\) is real and nonzero. By Proposition~\ref{prop:r-positivity-branch},
\[
\Im r_{j-1}(x_*)>0
\qquad
(j\ge 2),
\]
and for \(j=1\) the convention \(r_0=\infty\) gives \(\alpha_1=x_*+iy\). In either case,
\[
\Im \alpha_j(x_*)
=
y+\Im\!\left(-\frac{a}{r_{j-1}(x_*)}\right)
>0.
\]
Now Lemma~\ref{lem:angular-identity} gives
\[
\Im\!\bigl(\alpha_j(x_*)\,\overline{p_j(x_*)}\bigr)
=
-\Im\!\bigl(\beta_j(x_*)\,p_j(x_*)\bigr).
\]
Since \(p_j(x_*)\in\mathbb R\setminus\{0\}\), this simplifies to
\[
p_j(x_*)\,\Im \alpha_j(x_*)
=
-p_j(x_*)\,\Im \beta_j(x_*),
\]
hence
\[
\Im \beta_j(x_*)=-\Im \alpha_j(x_*)<0.
\]

On the other hand, the reflected ratio recursion gives
\[
\overline{\rho_j}
=
\frac{a}{\beta_j-s/p_j}.
\]
At \(x_*\), the denominator has the same imaginary part as \(\beta_j(x_*)\), namely a strictly negative one. Therefore
\[
\Im \overline{\rho_j}(x_*)>0,
\qquad\text{so}\qquad
\Im \rho_j(x_*)<0,
\]
contradicting the hypothesis \(\Im \rho_j>0\) on \(J\).

Thus \(\Im p_j\) never vanishes on \(J\). Since \(\Im p_j(x_0)>0\) at one point and \(J\) is connected, we conclude
\[
\Im p_j(x)>0
\qquad
(x\in J).
\]
\end{proof}

\begin{corollary}[Monotonicity on nonvanishing simple branches]\label{cor:monotonicity-on-branch}
Fix \(y>0\), and let \(J\) be as in Proposition~\ref{prop:r-positivity-branch}. Assume that at some \(x_0\in J\),
\[
\Im r_j(x_0)>0,
\qquad
j=1,\dots,n-1,
\]
and
\[
\Im p_j(x_0)>0,
\qquad
j=1,\dots,n.
\]
Then for every \(x\in J\),
\[
\frac{\partial}{\partial y}\bigl(s(x,y)^2\bigr)\ge 0.
\]
\end{corollary}

\begin{proof}
By Propositions~\ref{prop:r-positivity-branch} and \ref{prop:p-positivity-branch},
\[
\Im p_j(x)>0
\qquad
(j=1,\dots,n,\ x\in J).
\]
Using
\[
\Im p_j
=
-\frac{\Im(v_jv_{n+1-j})}{|v_j|^2},
\]
we obtain
\[
\Im(v_jv_{n+1-j})<0.
\]
Hence Proposition~\ref{prop:current-identity} gives
\[
J_j-aJ_{j-1}
=
y|v_j|^2+s(x,y)\,\Im(v_jv_{n+1-j})
\le y|v_j|^2.
\]
Inductively,
\[
J_j\le y\sum_{k=1}^j a^{\,j-k}|v_k|^2,
\]
and therefore
\[
(1-a)\sum_{j=1}^{n-1}J_j
\le
y\sum_{k=1}^{n-1}(1-a^{\,n-k})|v_k|^2
\le
y.
\]
Now Proposition~\ref{prop:current-identity} yields
\[
\frac{\partial}{\partial y}\bigl(s(x,y)^2\bigr)
=
2y-2(1-a)\sum_{j=1}^{n-1}J_j
\ge 0.
\]
\end{proof}

\begin{corollary}[Local continuation from the imaginary axis]\label{cor:local-continuation-from-x0}
Fix \(y_*>0\). Then there exists \(\delta>0\) such that
\[
\frac{\partial}{\partial y}\bigl(s(x,y)^2\bigr)\Big|_{y=y_*}\ge 0
\qquad
\text{for all }x\in(-\delta,\delta).
\]
\end{corollary}

\begin{proof}
By Proposition~\ref{prop:x=0-monotonicity}, at \(x=0\) one has a simple ground-state branch with
\[
\Im r_j(0,y_*)>0,
\qquad
j=1,\dots,n-1,
\]
and
\[
\Im p_j(0,y_*)>0,
\qquad
j=1,\dots,n.
\]
Since \(\lambda_1(0,y_*)\) is simple, standard analytic perturbation theory yields an open interval \(I\ni 0\) on which \(\lambda_1(x,y_*)\) remains simple and the corresponding singular vector is analytic in \(x\). The set
\[
\{x\in I:\ v_1(x)\cdots v_n(x)\neq 0\}
\]
is open and contains \(0\), so it contains some interval \((-\delta,\delta)\). On that interval, Corollary~\ref{cor:monotonicity-on-branch} applies and proves the claim.
\end{proof}

\begin{remark}[What remains open]\label{rem:remaining-gap-monotonicity}
The results above prove Conjecture~\ref{conj:vertical-monotonicity} in full on the slice \(x=0\), and they prove the differential inequality
\[
\frac{\partial}{\partial y}\bigl(s(x,y)^2\bigr)\ge 0
\]
on every simple branch in the \(x\)-variable on which the right singular vector has no zero entries. By Lemma~\ref{lem:boundary-entries-nonzero}, only interior zeros can occur. The remaining unresolved issue is therefore global continuation in \(x\): one must exclude, or control, possible interior zeros of the lowest singular vector and possible singular-value multiplicity points.
\end{remark}

\end{document}